%% Question bank chapter file - generated by bank_tools.py build (do not edit while the bank is still being built:
%% it is regenerated from the per-test data in Question Bank\data). Every question is one bankq block:
%%   \begin{bankq}[id=..., source={...}, paper=.., qno=.., type=..., subtopic={...}, answer={...}, answer_src=key|solved,
%%                 working={...}, difficulty=1-10, diff_parts={K S I C F}, diff_note={...},
%%                 group=<passage id or empty>, also_in={duplicate copies in other tests}]
%% difficulty = min(10, K+S+I+C+F-1): K concepts 1-3, S steps 1-4, I insight 0-2, C computation 0-2, F format 0-1
%%   ...question LaTeX, options, figures (\bfig / \bfigi -> img/)...
%%   \end{bankq}
%% Passages shared by several questions are bankpassage blocks; their questions carry group=<passage id>.
\documentclass[10pt]{article}
\usepackage[a4paper,margin=1.8cm]{geometry}
\usepackage[T1]{fontenc}
\usepackage{amsmath,amssymb,mathtools}
\usepackage{graphicx}
\usepackage[export]{adjustbox}
\usepackage{xcolor,needspace,keyval}
\usepackage{enumitem}
\usepackage{booktabs,array,tabularx}
\usepackage[hidelinks]{hyperref}
\setlength{\parindent}{0pt}
\setlength{\parskip}{3pt}
\makeatletter
\define@key{bq}{id}{\def\bq@id{#1}}\define@key{bq}{source}{\def\bq@src{#1}}\define@key{bq}{paper}{\def\bq@paper{#1}}
\define@key{bq}{qno}{\def\bq@qno{#1}}\define@key{bq}{type}{\def\bq@type{#1}}\define@key{bq}{subtopic}{\def\bq@sub{#1}}
\define@key{bq}{answer}{\def\bq@ans{#1}}\define@key{bq}{answer_src}{\def\bq@asrc{#1}}\define@key{bq}{working}{\def\bq@work{#1}}
\define@key{bq}{group}{\def\bq@group{#1}}\define@key{bq}{also_in}{\def\bq@also{#1}}
\define@key{bq}{difficulty}{\def\bq@diff{#1}}\define@key{bq}{diff_parts}{\def\bq@dparts{#1}}\define@key{bq}{diff_note}{\def\bq@dnote{#1}}
\newenvironment{bankq}[1][]{\def\bq@work{}\def\bq@also{}\def\bq@group{}\def\bq@diff{?}\def\bq@dparts{}\def\bq@dnote{}\setkeys{bq}{#1}%
  \par\needspace{4\baselineskip}\medskip\noindent{\small\textbf{\bq@id}\hfill\color{gray}[\bq@src, Paper \bq@paper, Q\bq@qno]
  \bq@type\ $\cdot$ \bq@sub\ $\cdot$ \textbf{D\bq@diff}}\par\nopagebreak\relax}
  {\par{\small\color{teal}Answer: \textbf{\bq@ans}\ifx\bq@work\@empty\else\ (solved: \bq@work)\fi
  \ifx\bq@also\@empty\else\par\color{gray}Also in: \bq@also\fi}\par}
\define@key{bp}{id}{\def\bp@id{#1}}\define@key{bp}{source}{\def\bp@src{#1}}\define@key{bp}{questions}{\def\bp@qs{#1}}
\newenvironment{bankpassage}[1][]{\setkeys{bp}{#1}\par\needspace{4\baselineskip}\bigskip\noindent
  {\small\textbf{Passage \bp@id}\hfill\color{gray}[\bp@src; questions \bp@qs]}\par\nopagebreak\relax\itshape}{\par}
\makeatother
\newcommand{\bfig}[2]{\par\smallskip{\centering\includegraphics[width=#1pt,max width=0.62\linewidth,max height=0.32\textheight]{img/#2}\par}\smallskip}
\newcommand{\bfigi}[2]{\includegraphics[width=#1pt,max width=\linewidth,max height=0.32\textheight,valign=c]{img/#2}}
\newcommand{\bmissing}[1]{\textit{[figure missing in the source conversion: \detokenize{#1}]}}
\newenvironment{qoptions}{\begin{itemize}[leftmargin=1.8em,itemsep=1pt,topsep=2pt,label={}]}{\end{itemize}}
\newcommand{\opt}[2]{\item[(#1)] #2}
\providecommand{\fwstart}{}\providecommand{\fwend}{}\providecommand{\qfullsplit}{}
\newcolumntype{L}[1]{>{\raggedright\arraybackslash}p{#1}}
\newcolumntype{C}[1]{>{\centering\arraybackslash}p{#1}}
\newcolumntype{R}[1]{>{\raggedleft\arraybackslash}p{#1}}
\providecommand{\sectionhdr}[1]{}\providecommand{\subsectionhdr}[1]{}


\begin{document}
{\Large\bfseries Isomerism --- Question Bank}\par\smallskip
{\small Chemistry chapter C15. Questions from all converted test papers, grouped by source series. Answers: printed key where it matches the paper, otherwise solved (answer\_src=solved, with working). Exact repeats across papers appear once, with also\_in listing the other copies.}\par\medskip\hrule

\section*{Score-I papers}
\begin{bankq}[id=SI01-P1-CHEM-S2-Q08, source={Score-I - Test 01 (16 Jun 2021)}, paper=1, qno=8, type=multi, subtopic={Optical}, answer={C,D}, answer_src=key, difficulty=5, diff_parts={K1 S2 I1 C1 F1}, diff_note={concepts: statements about two elimination reactions | steps: Hofmann versus Zaitsev; number of alkenes formed; stability; check each}]
Select correct statement for following reactions :

\textbf{Reaction : I}
\bfig{127.5}{SI01_p1_chem_s2_q08_1.png}
\textbf{Reaction : II}
\bfig{125.9}{SI01_p1_chem_s2_q08_2.png}
\begin{qoptions}
\opt{A}{Reaction I gives mixture of 2 enantiomers \& Reaction II gives mixture of two diastereomers.}
\opt{B}{Both reactions give mixture of 2 enantiomers.}
\opt{C}{Both reactions give mixture of 2 diastereomers.}
\opt{D}{Both reactions give 2 fraction of organic compounds.}
\end{qoptions}
\end{bankq}
\begin{bankq}[id=SI01-P1-CHEM-S3-Q04, source={Score-I - Test 01 (16 Jun 2021)}, paper=1, qno=4, type=numerical, subtopic={Structural}, answer={6}, answer_src=key, difficulty=7, diff_parts={K2 S2 I2 C1 F1}, diff_note={concepts: count, isomers on replacing one ligand in Ma3b3 isomers | steps: mer and fac isomers; replace one a by b; resulting Ma2b4 geometric forms; count | I: systematic enumeration of isomers from two parent isomers}]
(a) From Meridional and facial isomer of $[Ma_3b_3]^{n\pm}$ on replacement of only one `a' by `b', the number of isomer of the product obtained are respectively $x$ and $y$.

(b) The number of corner shared per tetrahedron for 2D-silicate is $z$.

Then $x+y+z$ is-
\end{bankq}
\begin{bankq}[id=SI01-P2-CHEM-S1-Q05, source={Score-I - Test 01 (16 Jun 2021)}, paper=2, qno=5, type=single, subtopic={Geometrical}, answer={C}, answer_src=key, difficulty=3, diff_parts={K1 S2 I1 C0 F0}, diff_note={concepts: stereochemistry of a drawn compound | steps: E Z assignment by priorities; chirality centres; check each}]
True statement(s) about this compound is/are:
\bfig{65.9}{SI01_p2_chem_s1_q05_1.png}
\begin{qoptions}
\opt{A}{It is E isomer}
\opt{B}{It is Z isomer}
\opt{C}{It is optically active}
\opt{D}{It has center of symmetry}
\end{qoptions}
\end{bankq}
\begin{bankq}[id=SI01-P2-CHEM-S2-Q12, source={Score-I - Test 01 (16 Jun 2021)}, paper=2, qno=12, type=multi, subtopic={Geometrical}, answer={A,B,D}, answer_src=key, difficulty=4, diff_parts={K1 S2 I1 C0 F1}, diff_note={concepts: indigotin properties | steps: trans form more stable; conjugation; geometrical isomerism; check each}]
\textbf{Correct} statements about indigotin is/are
\bfig{109.1}{SI01_p2_chem_s2_q12_1.png}
\begin{qoptions}
\opt{A}{It shows geometrical isomerism}
\opt{B}{Trans is more stable than cis}
\opt{C}{Cis is more stable than trans}
\opt{D}{It is a planar compound}
\end{qoptions}
\end{bankq}
\begin{bankq}[id=SI01-P2-CHEM-S2-Q13, source={Score-I - Test 01 (16 Jun 2021)}, paper=2, qno=13, type=multi, subtopic={Geometrical}, answer={A,B,C}, answer_src=key, difficulty=5, diff_parts={K1 S2 I1 C1 F1}, diff_note={concepts: statements about the drawn reaction | steps: identify product; geometrical isomerism; stereochemistry; check}]
\bfig{139.0}{SI01_p2_chem_s2_q12_2.png}
Which of the following statements is/are correct about the above reaction ?
\begin{qoptions}
\opt{A}{Product shows geometrical isomerism.}
\opt{B}{Product shows colour due to extended conjugation.}
\opt{C}{Electrophile attack at para position due to large size.}
\opt{D}{Electron withdrawing group in phenol increases rate of reaction.}
\end{qoptions}
\end{bankq}
\begin{bankq}[id=SI05-P1-CHEM-S3-Q04, source={Score-I - Test 05 (22 Jan 2022)}, paper=1, qno=4, type=numerical, subtopic={Optical}, answer={20}, answer_src=solved, working={Chiral monobromides of the five C6H14 isomers: hexane 2 constitutions, 2-methylpentane 3, 3-methylpentane 2 (one with two centres), 2,2-dimethylbutane 1, 2,3-dimethylbutane 1; counting enantiomers: 8 single-centre x 2 + 4 stereoisomers = 20. Key slot missing.}, difficulty=6, diff_parts={K2 S2 I1 C1 F1}, diff_note={concepts: optically active monobromides from isomers, C6H14 | steps: list five isomers of hexane; enumerate monobromination products; identify chiral ones; remove duplicates}]
How many optically active monobromides are possible on bromination of a mixture containing all isomers of C$_6$H$_{14}$
\end{bankq}
\begin{bankq}[id=SI06-P1-CHEM-S2-Q07, source={Score-I - Test 06 (28 Jan 2022)}, paper=1, qno=7, type=multi, subtopic={Geometrical}, answer={A,C,D}, answer_src=key, difficulty=4, diff_parts={K1 S1 I1 C1 F1}, diff_note={concepts: incorrect statements about E and Z isomers of chlorofluoropropene | steps: dipole moment, boiling point and solubility compare; check each}]
Which of the following incorrect for stereoisomers 2-chloro-1-fluoro propene.
\begin{qoptions}
\opt{A}{E--isomer is having more water solubility then Z--isomer.}
\opt{B}{Z--isomer is having more water solubility then E--isomer}
\opt{C}{E \& Z both isomers are equally water soluble}
\opt{D}{E--isomer is having intramolecular H--bonding.}
\end{qoptions}
\end{bankq}
\begin{bankq}[id=SI06-P1-CHEM-S2-Q08, source={Score-I - Test 06 (28 Jan 2022)}, paper=1, qno=8, type=multi, subtopic={Optical}, answer={C}, answer_src=key, difficulty=5, diff_parts={K1 S2 I1 C1 F1}, diff_note={concepts: R S configuration of a compound with two stereocentres | steps: assign priorities by CIP rules at each centre; determine R or S; match option}]
Correct R, S configuration for following compound is:
\bfig{116.1}{SI06_p1_chem_s2_q08_1.png}
\begin{qoptions}
\opt{A}{R at A, R at B}
\opt{B}{R at A, S at B}
\opt{C}{S at A, R at B}
\opt{D}{S at A, S at B}
\end{qoptions}
\end{bankq}
\begin{bankq}[id=SI06-P1-CHEM-S3-Q05, source={Score-I - Test 06 (28 Jan 2022)}, paper=1, qno=5, type=numerical, subtopic={Optical}, answer={3.00}, answer_src=key, difficulty=5, diff_parts={K1 S2 I1 C1 F1}, diff_note={concepts: count chiral compounds | steps: find compounds with no plane of symmetry; chiral centres; count}]
How many of the following compounds are chiral ?

(i) \bfigi{76.5}{SI06_p1_chem_s3_q05_i.png}

(ii) \bfigi{124.4}{SI06_p1_chem_s3_q05_ii.png}

(iii) \bfigi{52.4}{SI06_p1_chem_s3_q05_iii.png}

(iv) \bfigi{39.9}{SI06_p1_chem_s3_q05_iv.png}

(v) \bfigi{81.3}{SI06_p1_chem_s3_q05_v.png}

(vi) \bfigi{118.3}{SI06_p1_chem_s3_q05_vi.png}
\end{bankq}
\begin{bankq}[id=SI06-P1-CHEM-S3-Q06, source={Score-I - Test 06 (28 Jan 2022)}, paper=1, qno=6, type=numerical, subtopic={Optical}, answer={420.00}, answer_src=key, difficulty=6, diff_parts={K1 S3 I1 C1 F1}, diff_note={concepts: stereoisomers enantiomeric pairs and meso forms | steps: identify stereocentres and symmetry; total stereoisomers P; enantiomeric pairs Q; meso forms R; combine as PQR}]
\bfigi{126.3}{SI06_p1_chem_s3_q06_1.png}

For above compound's

P $\to$ total number of stereoisomers.

Q $\to$ total number of enantiomeric pairs

R $\to$ total number of meso compounds

Then give your answer in 3 digit's as PQR.
\end{bankq}
\begin{bankq}[id=SI09-P1-CHEM-S1-Q02, source={Score-I - Test 09 (16 Feb 2022)}, paper=1, qno=2, type=single, subtopic={Structural}, answer={B}, answer_src=key, difficulty=2, diff_parts={K1 S1 I1 C0 F0}, diff_note={concepts: most stable conformation of meso butane-2 3-diol | steps: intramolecular hydrogen bonding favours gauche OH conformer; choose}]
Among the following, the conformation that corresponds to the most stable conformation of \textit{meso}--butane--2,3--diol is --
\begin{qoptions}
\opt{A}{\bfigi{64.2}{SI09_p1_chem_s1_q02_A.png}}
\opt{B}{\bfigi{63.0}{SI09_p1_chem_s1_q02_B.png}}
\opt{C}{\bfigi{69.6}{SI09_p1_chem_s1_q02_C.png}}
\opt{D}{\bfigi{66.4}{SI09_p1_chem_s1_q02_D.png}}
\end{qoptions}
\end{bankq}
\begin{bankq}[id=SI09-P1-CHEM-S2-Q05, source={Score-I - Test 09 (16 Feb 2022)}, paper=1, qno=5, type=multi, subtopic={Optical}, answer={C,D}, answer_src=key, difficulty=6, diff_parts={K2 S1 I2 C1 F1}, diff_note={concepts: compounds giving a product, specific optical rotation on nitration | steps: identify chiral compounds whose nitro product is optically inactive; check each compound | I: nitration changes the stereo centre status}]
Given
\bfig{147.2}{SI09_p1_chem_s2_q05_1.png}
The compound(s), which on reaction with $\mathrm{HNO_3}$ will give the product having degree of rotation, $[\alpha]_D=-52.7^\circ$ is (are)
\begin{qoptions}
\opt{A}{\bfigi{48.3}{SI09_p1_chem_s2_q05_A.png}}
\opt{B}{\bfigi{48.7}{SI09_p1_chem_s2_q05_B.png}}
\opt{C}{\bfigi{49.2}{SI09_p1_chem_s2_q05_C.png}}
\opt{D}{\bfigi{49.5}{SI09_p1_chem_s2_q05_D.png}}
\end{qoptions}
\end{bankq}
\begin{bankq}[id=SI10-P1-CHEM-S3-Q02, source={Score-I - Test 10 (24 Feb 2022)}, paper=1, qno=2, type=matching, subtopic={Conformational isomerism}, answer={A$\to$R,T; B$\to$R,T; C$\to$Q,R; D$\to$Q,R,S}, answer_src=solved, working={Ethane and propane have only the staggered stable conformer. In 1,2-dichloroethane and 1-chloropropane the anti form is more stable than the gauche form; 1-chloropropane has the position isomer 2-chloropropane.}, difficulty=3, diff_parts={K1 S2 I0 C0 F1}, diff_note={concepts: conformations of alkanes and haloalkanes | steps: ethane and propane have only the staggered stable form; anti form is more stable in dichloroethane; anti form is more stable in chloropropane and it has the position isomer chloropropane}]
\begin{center}\footnotesize
\begin{tabular}{|p{0.55cm}|p{3.2cm}|p{0.6cm}|p{2.9cm}|}
\hline
\multicolumn{2}{|c|}{Column - I} & \multicolumn{2}{c|}{Column - II}\\ \hline
(A) & Ethane & (P) & Gauche form is more stable than anti form\\ \hline
(B) & Propane & (Q) & Anti form is more stable than gauche form\\ \hline
(C) & 1,2-Dichloroethane & (R) & Can show conformational isomerism\\ \hline
(D) & 1-Chloropropane & (S) & Can show position isomerism\\ \hline
 & & (T) & Only one stable conformation\\ \hline
\end{tabular}
\end{center}
\end{bankq}
\begin{bankq}[id=SI11-P1-CHEM-S1-Q04, source={Score-I - Test 11 (27 Feb 2022)}, paper=1, qno=4, type=single, subtopic={Optical}, answer={C}, answer_src=key, difficulty=4, diff_parts={K1 S2 I1 C1 F0}, diff_note={concepts: identical molecules among Newman projections | steps: convert each to a common conformation by rotating the rear carbon; compare configuration at each carbon; find the identical pair}]
Newman projections \textbf{P}, \textbf{Q}, \textbf{R} and \textbf{S} are shown below :
\bfig{275.9}{SI11_p1_chem_s1_q04_1.png}

Which one of the following options represents identical molecules ?
\begin{qoptions}
\opt{A}{\textbf{P} and \textbf{Q}}
\opt{B}{\textbf{Q} and \textbf{S}}
\opt{C}{\textbf{Q} and \textbf{R}}
\opt{D}{\textbf{R} and \textbf{S}}
\end{qoptions}
\end{bankq}
\begin{bankq}[id=SI11-P2-CHEM-S3-Q06, source={Score-I - Test 11 (27 Feb 2022)}, paper=2, qno=6, type=numerical, subtopic={Structural}, answer={6}, answer_src=key, difficulty=6, diff_parts={K2 S2 I1 C1 F1}, diff_note={concepts: number of isomers of C8H10O2 that are optically active, give colour with ferric chloride | steps: phenol with chiral centre; place OH on ring and chiral side chain; count positional and optical isomers; total}, also_in={JEE Advanced (2021 Series) - Test 02 (04 Feb 2021) P2 Q6}]
An organic compound ($\mathrm{C_8H_{10}O_2}$) rotates plane-polarized light. It produces pink color with neutral $\mathrm{FeCl_3}$ solution. What is the total number of all the possible isomers for this compound?
\end{bankq}
\begin{bankq}[id=SI17-P1-CHEM-S1-Q06, source={Score-I - Test 17 (17 Apr 2022)}, paper=1, qno=6, type=single, subtopic={Optical}, answer={A}, answer_src=key, difficulty=5, diff_parts={K1 S2 I2 C1 F0}, diff_note={concepts: percentages of enantiomers from specific rotation of a mixture | steps: optical purity = observed/pure = 3.5/13.9; excess of minus enantiomer; percentages (100 + e)/2 and (100 - e)/2 | I: convert optical purity to composition}]
(+)2--butanol has $[\theta]_\lambda^{25}=13.9^\circ$. A sample of 2-butanol containing both the enantiomers was found to have specific rotation value of $-3.5^\circ$ under similar condition. The percentage of the (+) and ($-$) enantiomer present in the sample are, respectively :
\begin{qoptions}
\opt{A}{37.4\% and 62.6\%}
\opt{B}{62.6\% and 37.4\%}
\opt{C}{42.2\% and 57.8\%}
\opt{D}{35.5\% and 64.5\%}
\end{qoptions}
\end{bankq}
\begin{bankq}[id=SI17-P2-CHEM-S3-Q06, source={Score-I - Test 17 (17 Apr 2022)}, paper=2, qno=6, type=numerical, subtopic={Geometrical}, answer={4}, answer_src=key, difficulty=5, diff_parts={K1 S2 I1 C1 F1}, diff_note={concepts: stereoisomer counts for a drawn compound | steps: identify chiral centres and double bonds; count optically active, meso and geometrical isomers; compute x + y - z}]
\bfigi{101.5}{SI17_p2_chem_s3_q06_1.png}\par
Number of optically active isomers = x

Number of meso compounds = y

Number of geometrical isomers = z

Report your answer as x + y $-$ z = ?
\end{bankq}
\begin{bankq}[id=SI18-P1-CHEM-S2-Q11, source={Score-I - Test 18 (24 Apr 2022)}, paper=1, qno=11, type=multi, subtopic={Optical}, answer={A,B,D}, answer_src=key, difficulty=3, diff_parts={K1 S1 I1 C0 F1}, diff_note={concepts: compound showing optical isomerism | steps: look for a carbon with four different groups or an axis of chirality; check}, also_in={JEE Advanced (2021 Series) - Test 01 (24 Jan 2021) P1 Q11}]
Which of the following compound can not exhibit optical isomerism ?
\begin{qoptions}
\opt{A}{\bfigi{59.0}{SI18_p1_chem_s2_q11_A.png}}
\opt{B}{\bfigi{50.8}{SI18_p1_chem_s2_q11_B.png}}
\opt{C}{\bfigi{32.3}{SI18_p1_chem_s2_q11_C.png}}
\opt{D}{CH$_3$--CH=C=CH$_2$}
\end{qoptions}
\end{bankq}
\begin{bankq}[id=SI22-P1-CHEM-S1-Q06, source={Score-I - Test 22 (04 Dec 2022)}, paper=1, qno=6, type=multi, subtopic={Optical}, answer={A,B,C,D}, answer_src=key, difficulty=5, diff_parts={K1 S2 I1 C1 F1}, diff_note={concepts: relationships among four stereoisomeric structures | steps: assign configuration at each stereocentre; mirror images are enantiomers; meso needs internal symmetry; diastereomers differ at some but not all centres}]
\bfigi{84.6}{SI22_p1_chem_s1_q06_1.png}\hspace{1.5em}\bfigi{88.2}{SI22_p1_chem_s1_q06_2.png}\par
\begin{qoptions}
\opt{A}{(i) and (iv) are enantiomers}
\opt{B}{(i) and (iii) are diastereomers}
\opt{C}{(ii) and (iii) are meso compounds}
\opt{D}{(ii) and (iii) are diastereomers.}
\end{qoptions}
\end{bankq}
\begin{bankpassage}[id=SI22-P1-CHEM-S2-G11-12, source={Score-I - Test 22 (04 Dec 2022)}, questions={11-12}]
List-I contains four isomers of C$_7$H$_{10}$O and List-II contains reactions \& properties.
\par\smallskip
\begin{center}\small
\begin{tabular}{|c|c|c|p{4.4cm}|}
\hline
 & List-I & & List-II\\
\hline
(I) & \bfigi{42.4}{SI22_p1_chem_s2_para11_12_1.png}\rule{0pt}{22pt} & (P) & Produces H$_2$ on reaction with metallic Na\\
\hline
(II) & \bfigi{27.7}{SI22_p1_chem_s2_para11_12_2.png}\rule{0pt}{22pt} & (Q) & Undergoes oxidation to aldehyde or ketone by PCC\\
\hline
(III) & \bfigi{38.0}{SI22_p1_chem_s2_para11_12_3.png}\rule{0pt}{22pt} & (R) & Cloudiness appears in Lucas Test at room temperature\\
\hline
(IV) & \bfigi{60.8}{SI22_p1_chem_s2_para11_12_4.png}\rule{0pt}{22pt} & (S) & Produces CH$_4$ with CH$_3$MgBr\\
\hline
 & & (T) & forms H--bond with H$_2$O\\
\hline
 & & (U) & Gives red color in Victor Meyer's test.\\
\hline
\end{tabular}
\end{center}
\end{bankpassage}
\begin{bankq}[id=SI22-P1-CHEM-S2-Q11, source={Score-I - Test 22 (04 Dec 2022)}, paper=1, qno=11, type=single, subtopic={Structural}, answer={A}, answer_src=key, difficulty=5, diff_parts={K1 S2 I2 C1 F0}, diff_note={concepts: reactions and properties of four isomers of C7H10O matched | steps: identify the four isomers as enone ring ketone alkynol or aldehyde; reactions with Br2 water Tollens iodoform and hydrogenation; decide which hold for isomers I and IV | I: structures deduced from reactivity of isomers}, group=SI22-P1-CHEM-S2-G11-12]
Which of the following options has the correct combination considering List-I and List-II.
\begin{qoptions}
\opt{A}{(I) -- P, Q, S, T, U}
\opt{B}{(I) -- P, Q, R, S, T}
\opt{C}{(IV) -- S, T, U}
\opt{D}{(IV) -- R, T, U}
\end{qoptions}
\end{bankq}
\begin{bankq}[id=SI22-P1-CHEM-S2-Q12, source={Score-I - Test 22 (04 Dec 2022)}, paper=1, qno=12, type=single, subtopic={Structural}, answer={C}, answer_src=key, difficulty=5, diff_parts={K1 S2 I2 C1 F0}, diff_note={concepts: properties of isomers II and III of C7H10O matched | steps: decide structure of isomers II and III; tests for unsaturation and carbonyl and acidity; match the properties P Q R S T U for each | I: more properties per isomer than in the previous question}, group=SI22-P1-CHEM-S2-G11-12]
Which of the following options has the correct combination considering List-I and List-II.
\begin{qoptions}
\opt{A}{(II) -- P, Q, R, T}
\opt{B}{(II) -- P, R, T, U}
\opt{C}{(III) -- P, Q, R, S, T}
\opt{D}{(III) -- Q, R, S, T, U}
\end{qoptions}
\end{bankq}
\begin{bankq}[id=SI22-P1-CHEM-S3-Q01, source={Score-I - Test 22 (04 Dec 2022)}, paper=1, qno=1, type=numerical, subtopic={Optical}, answer={8.00}, answer_src=key, difficulty=6, diff_parts={K2 S2 I1 C1 F1}, diff_note={concepts: number, stereoisomers of a polycyclic or polychiral compound | steps: count stereocentres and double bond geometries; account for symmetry reducing the total; 2^n minus meso forms}]
Calculate the number of stereoisomers for the following compound :
\bfig{223.0}{SI22_p1_chem_s3_q01_1.png}
\end{bankq}
\begin{bankq}[id=SI22-P1-CHEM-S3-Q03, source={Score-I - Test 22 (04 Dec 2022)}, paper=1, qno=3, type=numerical, subtopic={Geometrical}, answer={8.00}, answer_src=key, difficulty=5, diff_parts={K1 S2 I1 C1 F1}, diff_note={concepts: count molecules showing both geometrical and optical isomerism | steps: check allenes for chirality; polyenes; disubstituted cyclohexane with cis trans and chirality; count each}]
How many of the following molecules show geometrical as well as optical isomerism.\\
(i) pent--2,3--diene, \quad (ii) octa--2,3,6--trine, \quad (iii) 1,3--dimethyl cyclohexane
\bfig{222.4}{SI22_p1_chem_s3_q03_1.png}
\end{bankq}
\begin{bankq}[id=SI22-P1-CHEM-S3-Q06, source={Score-I - Test 22 (04 Dec 2022)}, paper=1, qno=6, type=numerical, subtopic={Optical}, answer={0.00}, answer_src=key, difficulty=6, diff_parts={K2 S2 I1 C1 F1}, diff_note={concepts: pairs of enantiomers, diastereomers from complete ozonolysis of a compound | steps: cleave each double bond to carbonyl fragments; identify chiral fragments and their stereoisomers; count pairs; add}]
Find the total number of pairs of enantiomers and pairs of diastereomers obtained upon complete ozonolysis of the following compound
\bfig{70.5}{SI22_p1_chem_s3_q06_1.png}
\end{bankq}
\begin{bankq}[id=SI31-P2-CHEM-S2-Q08, source={Score-I - Test 31 (29 Dec 2022)}, paper=2, qno=8, type=numerical, subtopic={Optical}, answer={18}, answer_src=key, difficulty=5, diff_parts={K1 S2 I1 C1 F1}, diff_note={concepts: number of possible products in the drawn reaction | steps: identify the reaction type; consider regioisomers and stereoisomers formed; count}]
Number of possible product formed
\bfig{88.2}{SI31_p2_chem_s2_q08_1.png}
\end{bankq}

\section*{Score-II papers}
\begin{bankq}[id=SII01-P1-CHEM-S3-Q05, source={Score-II - Test 01 (21 Mar 2021)}, paper=1, qno=5, type=numerical, subtopic={Optical}, answer={6.00}, answer_src=key, difficulty=5, diff_parts={K1 S2 I1 C1 F1}, diff_note={concepts: count optically active complexes | steps: check symmetry of each coordination compound; chelates cis forms optically active; count}]
How many of the following compounds can be optically active?
(a) [Co(en)$_2$ClBr] \quad (b) [Co(en)(NH$_3)_2$Cl$_2$]$^+$ \quad (c) [Cr(NH$_3)_4$Cl$_2$]$^+$ \quad (d) [Co(NH$_3)_3$Cl$_3$] \quad (e) [Co(en)$_3$]$^{3+}$ \quad (f) [Ni(dmg)$_2$] \quad (g) [Pb(EDTA)]$^{2-}$ \quad (h) [Cr(ox)$_2$(NH$_3)_2$]$^-$ \quad (i) [Cr(gly)$_3$]
\end{bankq}
\begin{bankq}[id=SII01-P2-CHEM-S3-Q07, source={Score-II - Test 01 (21 Mar 2021)}, paper=2, qno=7, type=numerical, subtopic={Structural}, answer={4}, answer_src=key, difficulty=5, diff_parts={K1 S2 I1 C1 F1}, diff_note={concepts: count compounds soluble in concentrated ammonia | steps: AgCl and some silver salts soluble via complex; others not; count}]
Amongst the following, the total number of compounds soluble in concentrated NH$_3$ solution is
\begin{center}
\begin{tabular}{@{}ll@{}}
(A) Ag$_2$CrO$_4$ & (B) Cu(OH)$_2\cdot$CuSO$_4$ \\
(C) PbSO$_3$ & (D) Al(OH)$_3$ \\
(E) Ni(OH)$_2$ & (F) Zn$_3$(PO$_4$)$_2$ \\
(G) BaSO$_4$ & (H) Bi(OH)$_2$NO$_3$ \\
(I) Mn(OH)$_2$ & \\
\end{tabular}
\end{center}
\end{bankq}
\begin{bankq}[id=SII02-P1-CHEM-S2-Q12, source={Score-II - Test 02 (28 Mar 2021)}, paper=1, qno=12, type=multi, subtopic={Optical}, answer={A,B,D}, answer_src=key, difficulty=4, diff_parts={K1 S2 I1 C0 F1}, diff_note={concepts: pairs of compounds differentiated by tests | steps: 2,4-DNP reacts with carbonyls; Tollens with aldehydes; iodoform with methyl ketones; check each pair}]
Select the incorrect options:

\begin{center}
(A)\; \raisebox{-0.4\height}{\bfigi{56.9}{SII02_q12_A1.png}} \; and \; \raisebox{-0.4\height}{\bfigi{56.9}{SII02_q12_A2.png}} \; can be differentiated by 2,4-DNP test.

\vspace{4pt}
(B)\; \raisebox{-0.4\height}{\bfigi{42.7}{SII02_q12_B1.png}} \; and \; \raisebox{-0.4\height}{\bfigi{42.7}{SII02_q12_B2.png}} \; can be differentiated by NaHCO$_3$ test.

\vspace{4pt}
(C)\; \raisebox{-0.4\height}{\bfigi{48.4}{SII02_q12_C.png}} $\xrightarrow[\text{Fuse}]{\text{NaOH}}$ Product give violet colour with neutral FeCl$_3$

\vspace{4pt}
(D) With NaOI, \; \raisebox{-0.4\height}{\bfigi{74.0}{SII02_q12_D.png}} \; can give yellow ppt.
\end{center}
\end{bankq}
\begin{bankq}[id=SII02-P1-CHEM-S3-Q06, source={Score-II - Test 02 (28 Mar 2021)}, paper=1, qno=6, type=numerical, subtopic={Optical}, answer={723.00}, answer_src=key, difficulty=6, diff_parts={K1 S3 I1 C1 F1}, diff_note={concepts: stereoisomers enantiomeric pairs and meso forms | steps: identify stereocentres and symmetry; total stereoisomers P; enantiomeric pairs Q; meso forms R; combine as PQR}]
For the compound shown below:

\begin{center}
\bfigi{184.6}{SII02_n6_alkene.png}
\end{center}

For above compound's\\
P $\to$ total number of stereoisomers.\\
Q $\to$ total number of enantiomeric pairs\\
R $\to$ total number of meso compounds\\
Then give your answer in 3 digit's as PQR
\end{bankq}
\begin{bankq}[id=SII03-P1-CHEM-S1-Q12, source={Score-II - Test 03 (05 May 2021)}, paper=1, qno=12, type=single, subtopic={Optical}, answer={C}, answer_src=key, difficulty=2, diff_parts={K1 S1 I1 C0 F0}, diff_note={concepts: compound showing optical isomerism | steps: look for a carbon with four different groups or an axis of chirality; check}]
Which of the following compound can exhibit optical isomerism?
\begin{qoptions}
\opt{A}{\bfigi{58.3}{SII03_p1_chem_s1_q12_2.png}}
\opt{B}{\bfigi{59.8}{SII03_p1_chem_s1_q12_1.png}}
\opt{C}{\bfigi{33.2}{SII03_p1_chem_s1_q12_3.png}}
\opt{D}{CH$_3$--CH=C=CH$_2$}
\end{qoptions}
\end{bankq}
\begin{bankq}[id=SII04-P1-CHEM-S1-Q03, source={Score-II - Test 04 (07 May 2021)}, paper=1, qno=3, type=single, subtopic={Optical}, answer={B}, answer_src=key, difficulty=2, diff_parts={K1 S2 I0 C0 F0}, diff_note={concepts: reactions giving a single optically active amine | steps: chiral centre and optical activity of the product; racemisation in the mechanism; stereospecific steps; count}]
Identify which reaction(s) produce(s) single optically active Amine molecule as product

\begin{center}
\bfigi{233.8}{SII04_p1chem_q3_amines.png}
\end{center}

\begin{qoptions}
\opt{A}{ii, iii, iv}
\opt{B}{iii only}
\opt{C}{iii, iv}
\opt{D}{none of these}
\end{qoptions}
\end{bankq}
\begin{bankq}[id=SII04-P2-CHEM-S1-Q01, source={Score-II - Test 04 (07 May 2021)}, paper=2, qno=1, type=multi, subtopic={Structural}, answer={A,B,D}, answer_src=key, difficulty=4, diff_parts={K1 S2 I1 C0 F1}, diff_note={concepts: stereoisomer relations among four isomers | steps: identify mirror images and non mirror images; enantiomer pairs; diastereomer pairs}]
\begin{center}
\bfigi{233.8}{SII04_p2chem_q1_fischer.png}
\end{center}

Which of the following statement are true about these isomers?

\begin{qoptions}
\opt{A}{I and II are a pair of enantiomer}
\opt{B}{III and IV are a pair of enantiomer}
\opt{C}{I and II are a pair of diastereomers}
\opt{D}{I and IV are a pair of diastereomers}
\end{qoptions}
\end{bankq}
\begin{bankq}[id=SII04-P2-CHEM-S4-Q03, source={Score-II - Test 04 (07 May 2021)}, paper=2, qno=3, type=numerical, subtopic={Structural}, answer={6}, answer_src=key, difficulty=6, diff_parts={K1 S2 I2 C1 F1}, diff_note={concepts: alcohols C7H14O giving a particular alkene on dehydration | steps: ring containing alcohols; rearrangements; major alkene by Saytzeff; count alcohols leading to the named alkene | I: ring size and rearrangement decide the alkene formed}]
How many alcohols belongs to molecular formula C$_7$H$_{14}$O (excluding stereoisomers) and having six or less carbon ring on dehydration with conc. H$_2$SO$_4$ at high temperature give

\begin{center}
\bfigi{44.3}{SII04_p2chem_q3_ring.png}
\end{center}

as major product.
\end{bankq}
\begin{bankq}[id=SII06-P1-CHEM-S2-Q12, source={Score-II - Test 06 (05 Jul 2022)}, paper=1, qno=12, type=multi, subtopic={Optical}, answer={B,C}, answer_src=key, difficulty=2, diff_parts={K1 S1 I0 C0 F1}, diff_note={concepts: resolvability of chiral molecules | steps: check for chiral centres or axes and for the plane of symmetry in each structure}]
Which of the following is/are Resolvable?
\begin{qoptions}
\opt{A}{\bfigi{66.3}{SII06_p1_chem_s2_q12_A.png}}
\opt{B}{\bfigi{71.5}{SII06_p1_chem_s2_q12_B.png}}
\opt{C}{\bfigi{68.9}{SII06_p1_chem_s2_q12_C.png}}
\opt{D}{\bfigi{71.5}{SII06_p1_chem_s2_q12_D.png}}
\end{qoptions}
\end{bankq}
\begin{bankq}[id=SII07-P1-CHEM-S2-Q01, source={Score-II - Test 07 (10 Jul 2022)}, paper=1, qno=1, type=numerical, subtopic={Structural}, answer={3}, answer_src=key, difficulty=3, diff_parts={K1 S1 I1 C0 F1}, diff_note={concepts: structural isomers of C2H2O | steps: ketene, ethynol and oxirene type structures; count distinct connectivities}]
How many structurally isomeric compounds can be formed using molecular formula C$_2$H$_2$O?
\end{bankq}
\begin{bankq}[id=SII13-P1-CHEM-S2-Q08, source={Score-II - Test 13 (06 Feb 2024)}, paper=1, qno=8, type=numerical, subtopic={Geometrical}, answer={4}, answer_src=key, difficulty=3, diff_parts={K1 S1 I1 C0 F1}, diff_note={concepts: geometrical isomers of a complex or alkene | steps: identify the double bonds or ligand positions that give cis trans isomers and count}]
How many Geometrical Isomers are possible for
\bfig{51.3}{SII13_p1_chem_s2_q08_1.png}
\end{bankq}
\begin{bankpassage}[id=SII14-P1-CHEM-S3-G5-6, source={Score-II - Test 14 (06 Feb 2024)}, questions={5-6}]
\bfig{316.2}{SII14_p1_chem_s3_q04_1.png}
\end{bankpassage}
\begin{bankq}[id=SII14-P1-CHEM-S3-Q05, source={Score-II - Test 14 (06 Feb 2024)}, paper=1, qno=5, type=numerical, subtopic={Optical}, answer={Bonus}, answer_src=solved, working={Printed key omits this question's answer and it was not solved; banked as Bonus.}, difficulty=5, diff_parts={K1 S2 I1 C1 F1}, diff_note={concepts: stereoisomers of the product in a named reaction sequence | steps: follow the sequence from the given reactant through each reagent; draw the product H; count stereocentres and double bond geometry; apply 2^n}, group=SII14-P1-CHEM-S3-G5-6]
How many stereoisomers are possible of product `H'?
\end{bankq}
\begin{bankq}[id=SII14-P1-CHEM-S3-Q06, source={Score-II - Test 14 (06 Feb 2024)}, paper=1, qno=6, type=numerical, subtopic={Optical}, answer={3.00}, answer_src=key, difficulty=5, diff_parts={K1 S2 I1 C1 F1}, diff_note={concepts: degree of unsaturation of the final product | steps: follow the sequence to the final product O; molecular formula; count the rings and pi bonds}, group=SII14-P1-CHEM-S3-G5-6]
Degree of unsaturation of product `O' is :
\end{bankq}
\begin{bankq}[id=SII19-P1-CHEM-S2-Q08, source={Score-II - Test 19 (20 Feb 2024)}, paper=1, qno=8, type=numerical, subtopic={Optical}, answer={8}, answer_src=key, difficulty=4, diff_parts={K1 S1 I1 C1 F1}, diff_note={concepts: chiral centres in fluorometholone | steps: identify each sp3 carbon with four different groups in the steroid; count 8}]
How many chiral centres are present in
\par\nobreak
\bfig{83.9}{SII19_p1_chem_s2_q08_1.png} Fluoromethanolone
\end{bankq}
\begin{bankq}[id=SII22-P1-CHEM-S2-Q10, source={Score-II - Test 22 (14 Mar 2024)}, paper=1, qno=10, type=numerical, subtopic={Geometrical}, answer={0}, answer_src=key, difficulty=3, diff_parts={K1 S1 I1 C0 F1}, diff_note={concepts: geometrical isomers of a polyene | steps: number of double bonds capable of cis-trans isomerism with unsymmetric ends; 2^n corrected for symmetry}]
\bfig{28.3}{SII22_p1_chem_s2_q10_1.png}
Total number of geometrical isomers of above compound are
\end{bankq}
\begin{bankq}[id=SII24-P1-CHEM-S1-Q14, source={Score-II - Test 24 (17 Mar 2024)}, paper=1, qno=14, type=single, subtopic={Optical}, answer={D}, answer_src=key, difficulty=3, diff_parts={K1 S1 I1 C1 F0}, diff_note={concepts: reaction giving a meso compound | steps: syn dihydroxylation with OsO4 of a cis alkene gives meso product; compare each alkene}]
Which of the following gives meso compound as product?
\begin{qoptions}
\opt{A}{$\mathrm{CH_3-CH=CH_2}\ \xrightarrow[\mathrm{NaHSO_3}]{\mathrm{OsO_4}}$}
\opt{B}{\bfigi{110.4}{SII24_p1_chem_s1_q14_1.png}}
\opt{C}{\bfigi{90.6}{SII24_p1_chem_s1_q14_2.png}}
\opt{D}{\bfigi{67.3}{SII24_p1_chem_s1_q14_3.png}}
\end{qoptions}
\end{bankq}
\begin{bankq}[id=SII25-P1-CHEM-S2-Q01, source={Score-II - Test 25 (20 Mar 2024)}, paper=1, qno=1, type=numerical, subtopic={Structural}, answer={6}, answer_src=key, difficulty=4, diff_parts={K1 S1 I1 C1 F1}, diff_note={concepts: double bond equivalents of a compound | steps: formula to DBE; count rings and pi bonds}]
Calculate double bond equivalent of compound (M)
\bfig{78.1}{SII25_p1_chem_s2_q01_1.png}
\end{bankq}
\begin{bankpassage}[id=SII25-P2-CHEM-S3-G1-2, source={Score-II - Test 25 (20 Mar 2024)}, questions={1-2}]
One of the best methods for the regioselective alkylation of ketones consists in an activation of the required $\alpha$-carbon atom. For example :
\bfig{310.6}{SII25_p2_chem_s3_para1_2_1.png}
\end{bankpassage}
\begin{bankq}[id=SII25-P2-CHEM-S3-Q01, source={Score-II - Test 25 (20 Mar 2024)}, paper=2, qno=1, type=numerical, subtopic={Structural}, answer={Bonus}, answer_src=solved, working={Printed key omits, truncates or corrupts this answer and the question was not solved; banked as Bonus.}, difficulty=4, diff_parts={K1 S1 I1 C1 F1}, diff_note={concepts: degree of unsaturation of compound D2 | steps: identify structure in a sequence; DBE}, group=SII25-P2-CHEM-S3-G1-2]
Degree of unsaturation of `D$_2$' is-
\end{bankq}
\begin{bankq}[id=SII25-P2-CHEM-S3-Q02, source={Score-II - Test 25 (20 Mar 2024)}, paper=2, qno=2, type=numerical, subtopic={Structural}, answer={Bonus}, answer_src=solved, working={Printed key omits, truncates or corrupts this answer and the question was not solved; banked as Bonus.}, difficulty=4, diff_parts={K1 S1 I1 C1 F1}, diff_note={concepts: number of carbons in compound D1 | steps: follow the sequence; count carbon}, group=SII25-P2-CHEM-S3-G1-2]
Number of carbon in `D$_1$'
\end{bankq}
\begin{bankq}[id=SII25-P2-CHEM-S4-Q07, source={Score-II - Test 25 (20 Mar 2024)}, paper=2, qno=7, type=numerical, subtopic={Optical}, answer={2}, answer_src=key, difficulty=4, diff_parts={K1 S1 I1 C1 F1}, diff_note={concepts: total stereoisomers of a compound | steps: count chiral centres and double bond geometry; subtract meso forms}]
Identify total number of stereoisomers of following compound :
\bfig{60.1}{SII25_p2_chem_s4_q07_1.png}
\end{bankq}
\begin{bankq}[id=SII26-P1-CHEM-S1-Q01, source={Score-II - Test 26 (20 Mar 2024)}, paper=1, qno=1, type=single, subtopic={Structural}, answer={C}, answer_src=key, difficulty=3, diff_parts={K1 S2 I0 C1 F0}, diff_note={concepts: counting isomers of ethers C5H12O and chlorides C4H9Cl | steps: count ether isomers by dividing the alkyl groups between oxygen; count structural isomers of butyl chloride including chain forms; N1 + N2}]
Total number of ether structure possible for molecular formula C$_5$H$_{12}$O : \fbox{$N_1$}\\
Total number of structure possible for molecular formula C$_4$H$_9$Cl : \fbox{$N_2$}\\
value of (N$_1$ + N$_2$) =
\begin{qoptions}
\opt{A}{6}
\opt{B}{8}
\opt{C}{10}
\opt{D}{9}
\end{qoptions}
\end{bankq}
\begin{bankq}[id=SII26-P1-CHEM-S2-Q05, source={Score-II - Test 26 (20 Mar 2024)}, paper=1, qno=5, type=multi, subtopic={Optical}, answer={A,B,D}, answer_src=key, difficulty=4, diff_parts={K1 S2 I1 C0 F1}, diff_note={concepts: configurational diastereomers among pairs of structures | steps: assign configurations R S or E Z in each pair; diastereomers differ at some but not all stereocentres; check four pairs}]
Which of following is/are configurational diastreomers :
\begin{qoptions}
\opt{A}{\bfigi{127.6}{SII26_p1_chem_s2_q05_A.png}}
\opt{B}{\bfigi{120.8}{SII26_p1_chem_s2_q05_B.png}}
\opt{C}{\bfigi{132.5}{SII26_p1_chem_s2_q05_C.png}}
\opt{D}{\bfigi{85.0}{SII26_p1_chem_s2_q05_4.png}}
\end{qoptions}
\end{bankq}
\begin{bankq}[id=SII26-P1-CHEM-S4-Q07, source={Score-II - Test 26 (20 Mar 2024)}, paper=1, qno=7, type=numerical, subtopic={Optical}, answer={4}, answer_src=key, difficulty=5, diff_parts={K1 S2 I1 C1 F1}, diff_note={concepts: count reactions giving racemic products | steps: prochiral carbonyls with Grignard give racemic alcohols; SN1 reactions racemise; addition to a planar centre; check each of six reactions}]
How many of following reactions will give racemic products :
\bfig{271.3}{SII26_p1_chem_s4_q07_1.png}

\bfig{219.4}{SII26_p1_chem_s4_q07_2.png}
\bfigi{52.7}{SII26_p1_chem_s4_q07_3.png}\ $\xrightarrow[\text{2. H}_2\text{O}]{\text{1. CH}_3\text{MgBr}}$
\bfig{74.9}{SII26_p1_chem_s4_q07_4.png}
\bfig{211.1}{SII26_p1_chem_s4_q07_5.png}
\bfig{128.6}{SII26_p1_chem_s4_q07_6.png}
\end{bankq}
\begin{bankq}[id=SII26-P1-CHEM-S4-Q08, source={Score-II - Test 26 (20 Mar 2024)}, paper=1, qno=8, type=numerical, subtopic={Structural}, answer={2}, answer_src=key, difficulty=4, diff_parts={K1 S2 I0 C1 F1}, diff_note={concepts: stable conformers with nonzero dipole moment | steps: identify staggered and gauche conformers of the compound; dipole of each by symmetry; count those with nonzero dipole}]
How many stable conformer of compound
\bfigi{47.0}{SII26_p1_chem_s4_q08_1.png}
possible which has nonzero dipole moment :
\end{bankq}
\begin{bankq}[id=SII29-P1-CHEM-S4-Q04, source={Score-II - Test 29 (26 Mar 2024)}, paper=1, qno=4, type=numerical, subtopic={Optical}, answer={3}, answer_src=key, difficulty=4, diff_parts={K1 S2 I1 C0 F1}, diff_note={concepts: counting optical stereoisomers with fixed and unfixed centres | steps: identify stereocentres and double bonds with fixed configuration; count those with unfixed configuration; combine allowing for symmetry}]
For the given compound [x]. The total number of optical active stereoisomers are -

\begin{center}
\bfigi{184.6}{SII29_p2chem_n4_diene.png}
\end{center}

Where --
This type of bond indicate that the configuration at the specific carbon and the geometry of the double bond is fixed.
But this type of bond indicate that the configuration at the specific carbon and the geometry of the double bond in not fixed.
\end{bankq}
\begin{bankq}[id=SII30-P1-CHEM-S2-Q07, source={Score-II - Test 30 (27 Mar 2024)}, paper=1, qno=7, type=numerical, subtopic={Optical}, answer={5}, answer_src=key, difficulty=4, diff_parts={K1 S1 I1 C1 F1}, diff_note={concepts: count compounds with chiral carbon | steps: check each structure for four different groups; count}]
Total number of compounds with Chiral carbon atoms from following is \underline{\hspace{2em}}.
\bfig{92.8}{SII30_p1_chem_s2_q07_1.png}
\end{bankq}
\begin{bankq}[id=SII31-P1-CHEM-S2-Q04, source={Score-II - Test 31 (29 Mar 2024)}, paper=1, qno=4, type=single, subtopic={Optical}, answer={3}, answer_src=key, difficulty=4, diff_parts={K1 S2 I1 C1 F0}, diff_note={concepts: count compounds not giving racemic mixture with Grignard | steps: carbonyl addition gives chiral centre racemic; symmetric ketones or aldehydes give achiral; count}]
How many of the following compounds don't give racemic mixture with CH$_3$MgBr :
\par\nobreak
\begin{center}\small
\begin{tabular}{|c|c|c|c|}\hline
(A) & \bfigi{50.2}{SII31_p1_chem_s2_q04_2.png} & (B) & \bfigi{43.2}{SII31_p1_chem_s2_q04_1.png}\\\hline
(C) & \bfigi{49.4}{SII31_p1_chem_s2_q04_4.png} & (D) & \bfigi{35.8}{SII31_p1_chem_s2_q04_3.png}\\\hline
\multicolumn{4}{|c|}{\bfigi{141.3}{SII31_p1_chem_s2_q04_5.png}}\\\hline
\end{tabular}
\end{center}
\end{bankq}
\begin{bankq}[id=SII31-P1-CHEM-S2-Q07, source={Score-II - Test 31 (29 Mar 2024)}, paper=1, qno=7, type=single, subtopic={Optical}, answer={3}, answer_src=key, difficulty=3, diff_parts={K1 S1 I1 C1 F0}, diff_note={concepts: count compounds not giving racemic mixture with CH3MgBr | steps: same logic as Grignard addition giving chiral product; count achiral products}]
How many of the following compounds don't give racemic mixture with CH$_3$MgBr :
\bfig{158.5}{SII31_p1_chem_s2_q07_1.png}
\end{bankq}
\begin{bankq}[id=SII33-P2-CHEM-S3-Q06, source={Score-II - Test 33 (01 Apr 2024)}, paper=2, qno=6, type=numerical, subtopic={Structural}, answer={1}, answer_src=key, difficulty=5, diff_parts={K1 S2 I1 C1 F1}, diff_note={concepts: reactants that cannot give the ozonolysis products | steps: from the given products work backwards; each candidate structure ozonolysed; compare; count the ones that cannot give the product}]
How many structures given below from a $-$ e that cannot be a reactant X of the following ozonolysis sequence?
\bfig{174.4}{SII33_p2_chem_s3_q06_1.png}
\bfig{231.9}{SII33_p2_chem_s3_q06_2.png}
\end{bankq}
\begin{bankq}[id=SII35-P1-CHEM-S3-Q08, source={Score-II - Test 35 (04 Feb 2025)}, paper=1, qno=8, type=numerical, subtopic={Structural}, answer={6}, answer_src=key, difficulty=5, diff_parts={K1 S2 I1 C1 F1}, diff_note={concepts: alcohols that can perform a given reaction | steps: molecular formula C15H16O degrees of unsaturation 8; structure types of secondary and tertiary alcohols that satisfy the reaction conditions; enumerate and count}]
\bfig{221.0}{SII35_p1_chem_s3_q08_1.png}
How many total 2$^\circ$ and 3$^\circ$ alcohol with molecular formula C$_{15}$H$_{16}$O can perform the above reactions.
\end{bankq}
\begin{bankq}[id=SII38-P1-CHEM-S3-Q02, source={Score-II - Test 38 (12 Feb 2025)}, paper=1, qno=2, type=single, subtopic={Structural}, answer={1.20}, answer_src=key, difficulty=6, diff_parts={K2 S3 I1 C1 F0}, diff_note={concepts: functional groups reduced by LAH, SBH and alcohol isomer count | steps: list functional groups in the compound; LAH reduces all carbonyl nitrile ester acid etc; SBH reduces aldehydes and ketones only; count isomeric alcohols of C5H12O with stereoisomers; combine}]
\bfig{74.9}{SII38_p1_chem_s3_q02_1.png}
In the above given compound\\
\textbf{X} = Total number of functional groups that can be reduced by LAH (Lithium aluminium hydride),\\
\textbf{Y} = Total number of functional groups that can be reduced by SBH (sodium brohydride)\\
\hspace*{1.5cm}\&\\
\textbf{Z} = Total number of isomeric alcohols of molecular formula C$_5$H$_{12}$O (including stereoisomers) which can give blue colour in \textbf{Victor - Meyer test}\\
Then value of $\left(\dfrac{X+Y}Z\right)$ = ...........
\end{bankq}
\begin{bankq}[id=SII39-P1-CHEM-S3-Q05, source={Score-II - Test 39 (02 Mar 2025)}, paper=1, qno=5, type=single, subtopic={Structural}, answer={6}, answer_src=key, difficulty=4, diff_parts={K1 S2 I1 C1 F0}, diff_note={concepts: number of enols of C4H8O | steps: list carbonyl isomers of C4H8O (butanal butanone and isobutanal); each tautomer enol including geometric isomers; count}]
The total number of enols possible for the formula C$_4$H$_8$O are
\end{bankq}
\begin{bankq}[id=SII42-P1-CHEM-S2-Q06, source={Score-II - Test 42 (19 Mar 2025)}, paper=1, qno=6, type=numerical, subtopic={Geometrical}, answer={2.00}, answer_src=key, difficulty=3, diff_parts={K1 S1 I1 C0 F1}, diff_note={concepts: geometrical isomerism of a product from a reaction scheme | steps: identify the product U from the reaction steps; count double bonds with cis and trans possibilities}]
Number of geometrical isomers for product (U) is

\begin{center}
\bfigi{233.8}{SII42_p2chem_n6_scheme.png}
\end{center}
\end{bankq}
\begin{bankq}[id=SII43-P1-CHEM-S1-Q06, source={Score-II - Test 43 (23 Mar 2025)}, paper=1, qno=6, type=single, subtopic={Geometrical}, answer={D}, answer_src=key, difficulty=3, diff_parts={K1 S1 I1 C1 F0}, diff_note={concepts: structure which cannot show geometrical isomerism | steps: check each alkene or ring for different groups on each end; identify the one lacking}]
\par\nobreak
\bfig{111.8}{SII43_p1_chem_s1_q06_1.png}
\begin{qoptions}
\opt{A}{C cannot exhibit GI}
\opt{B}{\bfigi{59.6}{SII43_p1_chem_s1_q06_2.png}}
\opt{C}{\bfigi{65.1}{SII43_p1_chem_s1_q06_4.png}}
\opt{D}{\bfigi{47.5}{SII43_p1_chem_s1_q06_3.png}}
\end{qoptions}
\end{bankq}
\begin{bankq}[id=SII45-P1-CHEM-S1-Q01, source={Score-II - Test 45 (26 Mar 2025)}, paper=1, qno=1, type=single, subtopic={Optical}, answer={B}, answer_src=key, difficulty=3, diff_parts={K1 S1 I1 C1 F0}, diff_note={concepts: number of optically active isomers in a reaction | steps: identify product with chiral centres; count stereoisomers excluding meso}]
The number of optically active isomers formed in the following reaction is
\bfig{95.1}{SII45_p1_chem_s1_q01_1.png}
\begin{qoptions}
\opt{A}{2}
\opt{B}{4}
\opt{C}{6}
\opt{D}{7}
\end{qoptions}
\end{bankq}
\begin{bankq}[id=SII47-P1-CHEM-S1-Q19, source={Score-II - Test 47 (30 Mar 2025)}, paper=1, qno=19, type=single, subtopic={Structural}, answer={C}, answer_src=key, difficulty=2, diff_parts={K1 S1 I1 C0 F0}, diff_note={concepts: stereochemical relationship between two molecules | steps: assign configurations or compare mirror images from the structures; decide identical enantiomers or diastereomers}]
What is the relationship between the molecules in the following pairs?
\bfig{216.9}{SII47_p1_chem_s1_q19_1.png}
\begin{qoptions}
\opt{A}{enantiomers}
\opt{B}{diastereomers}
\opt{C}{identical}
\opt{D}{structural isomers}
\end{qoptions}
\end{bankq}
\begin{bankq}[id=SII47-P1-CHEM-S1-Q20, source={Score-II - Test 47 (30 Mar 2025)}, paper=1, qno=20, type=single, subtopic={Optical}, answer={A}, answer_src=key, difficulty=2, diff_parts={K1 S1 I1 C0 F0}, diff_note={concepts: counting stereoisomers | steps: number of geometrical or optical isomers = 2^n with corrections; match each compound with the count}]
Match stereo isomers for above compounds:
\begin{center}
\resizebox{\linewidth}{!}{%
\begin{tabular}{|c|c|c|c|}
\hline
\multicolumn{2}{|c|}{\textbf{Column-I}} & \multicolumn{2}{c|}{\textbf{Column-II}} \\
\hline
(A) & $CH_3\text-CH{=}CH\text-CH{=}N\text-OH$ & (P) & 2 \\
\hline
(B) & \bfigi{69.2}{SII47_p1_chem_s1_q20_2.png} & (Q) & 4 \\
\hline
(C) & \bfigi{130.3}{SII47_p1_chem_s1_q20_3.png} & (R) & 3 \\
\hline
(D) & \begin{tabular}{@{}c@{}}$CH_3\text-CH{=}CH\text-CH{=}CH\text-CH\text-Ph$\\ $|$\\ $CH_3$\end{tabular} & (S) & 8 \\
\hline
\end{tabular}}
\end{center}
\begin{qoptions}
\opt{A}{A $\to$ Q, B $\to$ P, C $\to$ R, D $\to$ S}
\opt{B}{A $\to$ S, B $\to$ P, C $\to$ R, D $\to$ Q}
\opt{C}{A $\to$ Q, B $\to$ R, C $\to$ P, D $\to$ S}
\opt{D}{A $\to$ P, B $\to$ P, C $\to$ Q, D $\to$ S}
\end{qoptions}
\end{bankq}

\section*{Score Advanced papers}
\begin{bankq}[id=SA01-P1-CHEM-S2-Q10, source={Score Advanced - Test 01 (22 Sep 2020)}, paper=1, qno=10, type=multi, subtopic={Optical}, answer={A,B,C,D}, answer_src=key, difficulty=3, diff_parts={K1 S2 I0 C0 F1}, diff_note={concepts: decarboxylation of acids | steps: beta keto acid decarboxylates; second; third; check options}]
Consider the following compounds\par
I~\bfigi{46.9}{SA01_p1_chem_s2_q10_1.png}\ II~\bfigi{74.8}{SA01_p1_chem_s2_q10_2.png}\ III~\bfigi{75.1}{SA01_p1_chem_s2_q10_3.png}\par
Select the correct statements
\begin{qoptions}
\opt{A}{(I) undergoes decarboxylation faster than II or III on heating}
\opt{B}{(I) when heated in the presence of acetophenone produces a compound having tertiary alcoholic group}
\opt{C}{Compound (I) acts as bidentate ligand}
\opt{D}{The compound (III) has the highest melting point.}
\end{qoptions}
\end{bankq}
\begin{bankq}[id=SA01-P2-CHEM-S3-Q06, source={Score Advanced - Test 01 (22 Sep 2020)}, paper=2, qno=6, type=numerical, subtopic={Optical}, answer={3.00}, answer_src=key, difficulty=4, diff_parts={K1 S2 I1 C0 F1}, diff_note={concepts: optical isomers | steps: stereocentres; symmetry; total stereoisomers; optically active count}]
Calculate total number of optically active stereoisomers for following compound?\par
\bfig{177.6}{SA01_p2_chem_s3_q06_1.png}
\end{bankq}
\begin{bankq}[id=SA01-P2-CHEM-S3-Q07, source={Score Advanced - Test 01 (22 Sep 2020)}, paper=2, qno=7, type=numerical, subtopic={Optical}, answer={8.70}, answer_src=key, difficulty=5, diff_parts={K1 S2 I1 C1 F1}, diff_note={concepts: specific rotation calculation | steps: concentration in g/mL; path length in dm; observed rotation; specific rotation}, also_in={JEE Advanced (2021 Series) - Test 07 (25 Apr 2021) P2 Q7}]
Calculate the specific rotation of (+) $-$ glyceraldehyde, if a solution containing 2.0 gm/ 10ml is placed in a 1 gm polarimeter tube and its observed rotation at 25$^\circ$C using sodium D line is +1.74$^\circ$.
\end{bankq}
\begin{bankq}[id=SA02-P2-CHEM-S3-Q06, source={Score Advanced - Test 02 (13 May 2021)}, paper=2, qno=6, type=numerical, subtopic={Optical}, answer={6}, answer_src=solved, working={Eight stereoisomers from three stereogenic units; hexachlorocyclohexane has 2 optically active forms}, difficulty=7, diff_parts={K2 S2 I2 C1 F1}, diff_note={concepts: stereoisomers count, optical | steps: stereogenic units; hexachlorocyclohexane isomers; difference}, also_in={JEE Advanced (2021 Series) - Test 10 (13 May 2021) P2 Q6}]
(a) Total number of stereoisomers of given compound are x\par
\bfig{195.7}{SA02_p2_chem_s3_q06_1.png}
(b) Number of optically active stereoisomers of 1,2,3,4,5,6-hexachlorocyclohexane is y.\par
Then x $-$ y is
\end{bankq}
\begin{bankq}[id=SA04-P1-CHEM-S1-Q07, source={Score Advanced - Test 04 (04 Jul 2021)}, paper=1, qno=7, type=multi, subtopic={Optical}, answer={A,D}, answer_src=solved, working={A is 2S,3S,4R versus 2S,3R,4R (diastereomers). B mirror images. C Newman 2R,3R versus Fischer 2S,3S (enantiomers). D is cis meso drawn twice}, difficulty=8, diff_parts={K3 S1 I2 C2 F1}, diff_note={concepts: pairs that are not enantiomers, optical, assign CIP descriptors or compare mirror images in Fischer projections | steps: assign CIP descriptors or compare mirror images in Fischer projections}]
Which of the following pair of compounds doesn't represent enantiomeric pair :
\begin{qoptions}
\opt{A}{\bfigi{111.3}{SA04_p1_chem_s1_q07_1.png}}
\opt{B}{\bfigi{108.1}{SA04_p1_chem_s1_q07_2.png}}
\opt{C}{\bfigi{107.2}{SA04_p1_chem_s1_q07_3.png}}
\opt{D}{\bfigi{141.6}{SA04_p1_chem_s1_q07_4.png}}
\end{qoptions}
\end{bankq}
\begin{bankq}[id=SA04-P1-CHEM-S1-Q08, source={Score Advanced - Test 04 (04 Jul 2021)}, paper=1, qno=8, type=multi, subtopic={Optical}, answer={A,B,C}, answer_src=solved, working={A cis and trans dichlorocyclobutane, B different diols (Et versus Me), C diastereomeric aldoses can be separated; D (R and S) cannot}, difficulty=6, diff_parts={K2 S1 I2 C1 F1}, diff_note={concepts: separation of stereoisomers by distillation, optical | steps: diastereomers and constitutional isomers have different boiling points; enantiomers do not}]
Identify pair(s) of compounds which can be seperate by fractional distillation ?
\begin{qoptions}
\opt{A}{\bfigi{148.0}{SA04_p1_chem_s1_q08_1.png}}
\opt{B}{\bfigi{106.0}{SA04_p1_chem_s1_q08_2.png}}
\opt{C}{\bfigi{105.7}{SA04_p1_chem_s1_q08_3.png}}
\opt{D}{\bfigi{96.7}{SA04_p1_chem_s1_q08_4.png}}
\end{qoptions}
\end{bankq}
\begin{bankq}[id=SA06-P1-CHEM-S1-Q04, source={Score Advanced - Test 06 (04 Sep 2021)}, paper=1, qno=4, type=multi, subtopic={Optical}, answer={B,C,D}, answer_src=key, difficulty=4, diff_parts={K1 S2 I1 C0 F1}, diff_note={concepts: reductive ozonolysis and optical activity | steps: break the double bond; the carbonyl products; chiral centres in the product; check each reactant}]
\par
\bfig{208.2}{SA06_p1_chem_s1_q04_1.png}
Which of the reactant will formed optically active product after reductive ozonolysis.
\begin{qoptions}
\opt{A}{I}
\opt{B}{II}
\opt{C}{III}
\opt{D}{IV}
\end{qoptions}
\end{bankq}
\begin{bankq}[id=SA06-P1-CHEM-S3-Q03, source={Score Advanced - Test 06 (04 Sep 2021)}, paper=1, qno=3, type=numerical, subtopic={Optical}, answer={2}, answer_src=key, difficulty=5, diff_parts={K1 S2 I1 C1 F1}, diff_note={concepts: optically pure carbonyl compounds forming chiral alcohols | steps: list isomers of C5H10O with a stereocentre; reduction creates a new centre or retains; count}]
x = number of optically pure aldehyde or ketone having molecular formula of C$_5$H$_{10}$O give chiral alcohol but not racemic mixture with NaBH$_4$ / CH$_3$--OH. (Consider enantiomer separately).
\end{bankq}
\begin{bankq}[id=SA08-P1-CHEM-S3-Q03, source={Score Advanced - Test 08 (22 Apr 2023)}, paper=1, qno=3, type=numerical, subtopic={Optical}, answer={0}, answer_src=key, difficulty=5, diff_parts={K1 S2 I1 C1 F1}, diff_note={concepts: optically active products of exhaustive bromination | steps: count positional isomers; chirality; stereoisomers; compute}]
\bfigi{110.2}{SA08_p1_chem_s3_q03_1.png}\par
How many optically active compounds having molecular formula C$_8$H$_8$Br$_8$ will be obtained in above reaction ?
\end{bankq}
\begin{bankq}[id=SA08-P2-CHEM-S3-Q05, source={Score Advanced - Test 08 (22 Apr 2023)}, paper=2, qno=5, type=numerical, subtopic={Optical}, answer={5}, answer_src=key, difficulty=2, diff_parts={K1 S1 I0 C0 F1}, diff_note={concepts: optical activity count | steps: chirality in each structure; count}]
Find the total number of optically active compounds\par
\bfigi{69.8}{SA08_p2_chem_s3_q05_2.png} \quad $\text{CH}_3-\text{CH}=\text{C}=\text{CH}-\text{CH}_3$, \quad \bfigi{120.8}{SA08_p2_chem_s3_q05_1.png}\par
\bfigi{198.7}{SA08_p2_chem_s3_q05_3.png}
\end{bankq}
\begin{bankq}[id=SA12-P1-CHEM-S3-Q01, source={Score Advanced - Test 12 (11 May 2023)}, paper=1, qno=1, type=numerical, subtopic={Structural}, answer={3}, answer_src=key, difficulty=3, diff_parts={K1 S1 I1 C0 F1}, diff_note={concepts: conformers with dipole moment | steps: staggered conformers; dipole of each}, also_in={Score Advanced - Test 29 (01 May 2024) P2 Q1}]
The total number(s) of stable conformers with non-zero dipole moment for the following compound is/are\par
\bfig{47.7}{SA12_p1_chem_s3_q01_1.png}
\end{bankq}
\begin{bankq}[id=SA13-P1-CHEM-S2-Q08, source={Score Advanced - Test 13 (13 May 2023)}, paper=1, qno=8, type=single, subtopic={Optical}, answer={5}, answer_src=key, difficulty=2, diff_parts={K1 S1 I1 C0 F0}, diff_note={concepts: optical activity with C2 axis | steps: chirality of each; presence of C2}]
Identify number of molecules which have C$_2$-axis of symmetry and they are optically active.\par
cis[Pt(NH$_3$)$_2$Cl$_2$] , 1,3dibromoallene ,\par
cis[Co(en)$_2$Cl$_2$]Cl , trans[Co(en)$_2$Br$_2$]NO$_3$ ,\par
trans 1-2dichlorocyclopentane ,\par
Na$_3$[Fe(C$_2$O$_4$)$_3$] , K[Co(EDTA)]\par
cis[Pt(NH$_3$)$_4$Cl$_2$]Br$_2$\par
\end{bankq}
\begin{bankq}[id=SA15-P1-CHEM-S1-Q03, source={Score Advanced - Test 15 (18 May 2023)}, paper=1, qno=3, type=single, subtopic={Structural}, answer={A}, answer_src=solved, working={Only IV}, difficulty=4, diff_parts={K1 S2 I1 C1 F0}, diff_note={concepts: most stable conformer of substituted molecules being achiral | steps: draw anti and gauche conformers; check symmetry of the lowest energy conformer for each compound; chirality}]
In which of the following cases, the most stable conformer is achiral?\par
\bfig{80.1}{SA15_p1_chem_s1_q03_1.png}
\begin{qoptions}
\opt{A}{Only IV}
\opt{B}{I and IV}
\opt{C}{I, II and III}
\opt{D}{In all the cases}
\end{qoptions}
\end{bankq}
\begin{bankq}[id=SA15-P1-CHEM-S3-Q03, source={Score Advanced - Test 15 (18 May 2023)}, paper=1, qno=3, type=numerical, subtopic={Optical}, answer={3}, answer_src=key, difficulty=5, diff_parts={K1 S2 I1 C1 F1}, diff_note={concepts: count correct stereochemical statements | steps: identify products X Y M N A B from the reactions; relation between each pair enantiomer diastereomer or identical; evaluate four statements}]
In the given reactions:-\par
\bfig{163.8}{SA15_p1_chem_s3_q03_1.png}
Number of correct statement(s) among the following are :-\par
(1) X and Y are enantiomers\par
(2) M and N are enantiomers\par
(3) A and B are diastereomers\par
(4) (X and Y) are identical with (A and B)\par
(5) (P and Q) are identical with (M and N)\par
(6) M and N are geometrical isomers
\end{bankq}
\begin{bankq}[id=SA16-P2-CHEM-S1-Q03, source={Score Advanced - Test 16 (20 May 2023)}, paper=2, qno=3, type=multi, subtopic={Structural}, answer={C,D}, answer_src=key, difficulty=5, diff_parts={K1 S2 I1 C1 F1}, diff_note={concepts: spectral or tests data to identify compounds | steps: analyze given data columns for compounds X and others; match tests; choose the correct statements}]
Analyze the following data and choose the \textbf{CORRECT} option(s)\par
{\footnotesize\setlength{\tabcolsep}{2pt}\noindent\begin{tabular}{|L{3.9cm}|L{3.9cm}|}\hline
\centering\textbf{Compound [X]} & \centering\textbf{Compound [Y]}\tabularnewline\hline
M.F : C$_4$H$_6$O (acyclic compound) & M.F : C$_4$H$_6$O (acyclic compound)\\\hline
Gives positive tollens test & Gives positive tollens test\\\hline
Gives positive bromine water test & Gives positive bromine water test\\\hline
Can show tautomerism & Cannot show tautomerism\\\hline
\end{tabular}}\par
\begin{qoptions}
\opt{A}{\textbf{[X]} \& \textbf{[Y]} are positional isomers}
\opt{B}{\textbf{[Y]} is an aldol condensation product of acetaldehyde}
\opt{C}{\textbf{[X]} on treatment with $\Delta_3$O$^+$, number of replaceable hydrogens with Deuterium are 4}
\opt{D}{\textbf{[X]} \& \textbf{[Y]} are chain isomers}
\end{qoptions}
\end{bankq}
\begin{bankq}[id=SA18-P1-CHEM-S3-Q02, source={Score Advanced - Test 18 (27 May 2023)}, paper=1, qno=2, type=numerical, subtopic={Structural}, answer={391.30}, answer_src=key, difficulty=5, diff_parts={K1 S2 I1 C1 F1}, diff_note={concepts: isomer count for C4H8O | steps: open chain carbonyl compounds 3; unsaturated alcohols count; ethers aliphatic and alicyclic; count each category and combine}]
C$_4$H$_8$O will have $x$ open chain carbonyl compounds, $y$ unsaturated alcohols and $z$ ethers ($z_1$ - aliphatic, $z_2$ - alicyclic). What is $[100x+10y+z-0.1\,(z_1-z_2)]$ ?
\end{bankq}
\begin{bankq}[id=SA19-P1-CHEM-S2-Q10, source={Score Advanced - Test 19 (04 Apr 2024)}, paper=1, qno=10, type=multi, subtopic={Optical}, answer={B,D}, answer_src=key, difficulty=5, diff_parts={K1 S2 I1 C1 F1}, diff_note={concepts: major product of the drawn rearrangement or addition | steps: identify mechanism for substrate; apply carbocation rearrangement or regiochemistry; choose}]
\par
\bfig{232.5}{SA19_p1_chem_s2_q10_1.png}
\begin{qoptions}
\opt{A}{\bfigi{72.2}{SA19_p1_chem_s2_q10_A.png}}
\opt{B}{\bfigi{115.1}{SA19_p1_chem_s2_q10_B.png}}
\opt{C}{\bfigi{87.5}{SA19_p1_chem_s2_q10_C.png}}
\opt{D}{\bfigi{116.3}{SA19_p1_chem_s2_q10_D.png}}
\end{qoptions}
\end{bankq}
\begin{bankq}[id=SA22-P2-CHEM-S2-Q08, source={Score Advanced - Test 22 (10 Apr 2024)}, paper=2, qno=8, type=multi, subtopic={Structural}, answer={A,B,C,D}, answer_src=key, difficulty=4, diff_parts={K1 S1 I1 C1 F1}, diff_note={concepts: isomerisation of a drawn compound standing in pentane | steps: a Cope or Claisen rearrangement or valence tautomerism converts to constitutional isomer; check statements}]
When the compounds of the following type are allowed to stand in pentane, they are converted into a constitutional isomer. Which of the following is/are correct about the constitutional isomers (A) and (B)?\par
\bfig{235.8}{SA22_p2_chem_s2_q08_1.png}
\begin{qoptions}
\opt{A}{(A) \& (B) are functional isomers}
\opt{B}{Both give same products on hydrolysis}
\opt{C}{Hydrolysis products can be separated by dil.HCl}
\opt{D}{Hydrolysis products can be separated by dil.NaOH}
\end{qoptions}
\end{bankq}
\begin{bankq}[id=SA22-P2-CHEM-S2-Q10, source={Score Advanced - Test 22 (10 Apr 2024)}, paper=2, qno=10, type=multi, subtopic={Optical}, answer={A,B,C}, answer_src=key, difficulty=5, diff_parts={K1 S2 I1 C1 F1}, diff_note={concepts: relative hydrogenation rates and heats for two alkenes | steps: more substituted alkene is more stable and hydrogenates slower; heat of hydrogenation smaller for the more stable one; check statements}]
Which of the following statement/s is/are correct about compounds (A) and (B)?\par
\bfig{167.3}{SA22_p2_chem_s2_q10_1.png}
\begin{qoptions}
\opt{A}{The rate of hydrogenation of (A) is more than (B) under similar conditions}
\opt{B}{The compounds (A) and (B) are stereo isomers.}
\opt{C}{Catalytic hydrogenation of (A) or (B) gives two diastereomers in unequal ratio.}
\opt{D}{The pair of diastereomers can be resolved in to optically active compounds.}
\end{qoptions}
\end{bankq}
\begin{bankq}[id=SA23-P1-CHEM-S1-Q02, source={Score Advanced - Test 23 (14 Apr 2024)}, paper=1, qno=2, type=single, subtopic={Optical}, answer={A}, answer_src=key, difficulty=4, diff_parts={K1 S2 I1 C1 F0}, diff_note={concepts: reactions giving diastereomers as major product | steps: addition of bromine to alkenes gives anti addition; hydrogenation syn; formation of two stereocentres leading to diastereomers; check each reaction}]
Identify the reaction(s) which yield diastereomers as a major product?
\begin{qoptions}
\opt{A}{\bfigi{101.1}{SA23_p1_chem_s1_q02_A.png}}
\opt{B}{\bfigi{114.4}{SA23_p1_chem_s1_q02_B.png}}
\opt{C}{\bfigi{131.1}{SA23_p1_chem_s1_q02_C.png}}
\opt{D}{\bfigi{135.2}{SA23_p1_chem_s1_q02_D.png}}
\end{qoptions}
\end{bankq}
\begin{bankq}[id=SA23-P1-CHEM-S2-Q08, source={Score Advanced - Test 23 (14 Apr 2024)}, paper=1, qno=8, type=multi, subtopic={Geometrical}, answer={B,C}, answer_src=key, difficulty=5, diff_parts={K1 S2 I1 C1 F1}, diff_note={concepts: distinguishing and properties of drawn compounds | steps: tests that differentiate P and Q; stereochemistry; resonance; check statements}]
\bfig{152.0}{SA23_p1_chem_s2_q08_1.png}
Considering above compounds, identify the correct statement(s)?
\begin{qoptions}
\opt{A}{Compound (\textbf{P}) and Compound (\textbf{Q}) can be differentiated by NaHCO$_3$ as well as aq.NaOH}
\opt{B}{Compound (\textbf{Q}) and Compound (\textbf{R}) can be differentiated by Neutral FeCl$_3$}
\opt{C}{Compound (\textbf{R}) and Compound (\textbf{S}) can be differentiated by NaHSO$_3$}
\opt{D}{Only compound (\textbf{R}) gives Mustard oil smell with CHCl$_3$ -KOH}
\end{qoptions}
\end{bankq}
\begin{bankq}[id=SA23-P2-CHEM-S4-Q09, source={Score Advanced - Test 23 (14 Apr 2024)}, paper=2, qno=9, type=numerical, subtopic={Optical}, answer={6}, answer_src=key, difficulty=5, diff_parts={K1 S2 I1 C1 F1}, diff_note={concepts: optically active stereoisomers of the given compound | steps: count stereocentres and geometric elements; account for meso forms; number of optically active isomers}]
For the given compound, the total number of optically active stereoisomers is?\par
\bfig{316.2}{SA23_p2_chem_s4_q09_1.png}
\end{bankq}
\begin{bankq}[id=SA24-P2-CHEM-S3-Q07, source={Score Advanced - Test 24 (17 Apr 2024)}, paper=2, qno=7, type=numerical, subtopic={Geometrical}, answer={4}, answer_src=key, difficulty=5, diff_parts={K1 S2 I1 C1 F1}, diff_note={concepts: geometrical isomers of 1 2 4-tribromocyclohexane | steps: three stereocentres on a ring; cis trans relationships: four geometrical isomers; count}]
Total no. of geometrical isomers possible for 1, 2, 4-tribromocyclohexane are?
\end{bankq}
\begin{bankq}[id=SA24-P2-CHEM-S3-Q08, source={Score Advanced - Test 24 (17 Apr 2024)}, paper=2, qno=8, type=numerical, subtopic={Optical}, answer={3}, answer_src=key, difficulty=5, diff_parts={K1 S2 I1 C1 F1}, diff_note={concepts: stereoisomers of 1 2-dimethylcyclopropane | steps: cis form is meso; trans form exists as an enantiomeric pair; total 3 stereoisomers}]
Total no. of stereoisomers possible for 1, 2-dimethylcyclopropane are?
\end{bankq}
\begin{bankq}[id=SA25-P1-CHEM-S4-Q06, source={Score Advanced - Test 25 (21 Apr 2024)}, paper=1, qno=6, type=numerical, subtopic={Optical}, answer={2}, answer_src=key, difficulty=5, diff_parts={K1 S2 I1 C1 F1}, diff_note={concepts: count products from reduction of all isomeric ketones of MW 86 | steps: ketones C5H10O: 2-pentanone 3-pentanone 3-methyl-2-butanone; reduction with NaBH4 gives alcohols; stereoisomers (2-pentanol racemic 3-methyl-2-butanol); count distinct products}, also_in={Score Advanced - Test 26 (21 Apr 2024) P1 Q6}]
Consider all possible isomeric ketones including stereoisomers of MW = 86, All these isomers are independently reacted with NaBH$_4$ (NOTE : stereoisomers are also reacted separately).\par
The total number of ketones that give a racemic product(s) is/are.
\end{bankq}
\begin{bankq}[id=SA25-P2-CHEM-S4-Q09, source={Score Advanced - Test 25 (21 Apr 2024)}, paper=2, qno=9, type=single, subtopic={Optical}, answer={7}, answer_src=key, difficulty=4, diff_parts={K1 S2 I1 C1 F0}, diff_note={concepts: reagents giving a redox reaction with solution A in acid | steps: identify A from the figure; test each listed reagent for redox behaviour; count}, also_in={Score Advanced - Test 26 (21 Apr 2024) P2 Q9}]
\bfig{190.3}{SA25_p2_chem_s4_q09_F.png}
Total number of reagent/s which give redox reaction with solution (A) in acidic medium among the following.\par
Pb(CH$_3$COO)$_2$ , Ba(NO$_3$)$_2$ , NO$_2^-$ , Sn$^{2+}$ , SO$_2$ , Fe$^{2+}$ , H$_2$S , I$^-$ , C$_2$H$_5$OH
\end{bankq}
\begin{bankq}[id=SA27-P1-CHEM-S1-Q02, source={Score Advanced - Test 27 (24 Apr 2024)}, paper=1, qno=2, type=single, subtopic={Structural}, answer={C}, answer_src=key, difficulty=4, diff_parts={K1 S2 I1 C1 F0}, diff_note={concepts: structural isomer of a given structure | steps: molecular formula of the given compound; compare formulas of the options; same formula different connectivity; select}]
The structural isomer of given structure is\par
\bfig{42.1}{SA27_p1_chem_s1_q02_1.png}
\begin{qoptions}
\opt{A}{\bfigi{31.3}{SA27_p1_chem_s1_q02_A.png}}
\opt{B}{\bfigi{46.5}{SA27_p1_chem_s1_q02_B.png}}
\opt{C}{\bfigi{33.8}{SA27_p1_chem_s1_q02_C.png}}
\opt{D}{\bfigi{33.7}{SA27_p1_chem_s1_q02_D.png}}
\end{qoptions}
\end{bankq}
\begin{bankq}[id=SA27-P2-CHEM-S1-Q02, source={Score Advanced - Test 27 (24 Apr 2024)}, paper=2, qno=2, type=multi, subtopic={Structural}, answer={B}, answer_src=solved, working={s-cis and s-trans butadiene (C) are conformers and the rotated hexadiene forms are conformers, but the (Z,Z)-diene versus the (E,E)-diene in B differ in configuration (cis/trans isomers) so B is not a conformer pair.}, difficulty=5, diff_parts={K1 S2 I1 C1 F1}, diff_note={concepts: which pair is not conformers | steps: conformers interconvert by rotation about single bonds; Newman or sawhorse pairs; check whether the two drawings are same molecule rotated or configurational isomers}]
Which of the following in not a correct conformer pair?
\begin{qoptions}
\opt{A}{\bfigi{105.8}{SA27_p2_chem_s1_q02_A.png}}
\opt{B}{\bfigi{101.1}{SA27_p2_chem_s1_q02_B.png}}
\opt{C}{\bfigi{81.1}{SA27_p2_chem_s1_q02_C.png}}
\opt{D}{\bfigi{119.0}{SA27_p2_chem_s1_q02_D.png}}
\end{qoptions}
\end{bankq}
\begin{bankq}[id=SA28-P1-CHEM-S3-Q07, source={Score Advanced - Test 28 (27 Apr 2024)}, paper=1, qno=7, type=numerical, subtopic={Optical}, answer={9}, answer_src=key, difficulty=5, diff_parts={K1 S2 I1 C1 F1}, diff_note={concepts: total stereoisomers of three coordination compounds | steps: tris-oxalato chromate: 2 (enantiomers); cis and trans PtCl2(en)2 : cis chiral 2 plus trans 1 = 3; third complex Cr(NH3)2Cl2 count; sum x y z}]
Find the total stereoisomers for :\par
(a) [Cr(C$_2$O$_4$)$_3$]$^{3-}$ = x\par
(b) [PtCl$_2$(en)$_2$]$^{2+}$ = y\par
(c) [Cr(NH$_3$)$_2$Cl$_2$ (en)]$^+$ = z\par
Hence write the value of (x + y + z).
\end{bankq}
\begin{bankq}[id=SA28-P2-CHEM-S2-Q07, source={Score Advanced - Test 28 (27 Apr 2024)}, paper=2, qno=7, type=single, subtopic={Optical}, answer={A}, answer_src=key, difficulty=4, diff_parts={K1 S2 I1 C1 F0}, diff_note={concepts: matching reactions  kinetic order or properties | steps: order from rate law units; half life dependence; examples; match}]
Matching List.\par
{\footnotesize\setlength{\tabcolsep}{1.5pt}\noindent\begin{tabular}{|c|L{3.9cm}|c|L{2.4cm}|}
\hline
 & \multicolumn{1}{>{\centering\arraybackslash}p{3.9cm}|}{List-I} & & \multicolumn{1}{>{\centering\arraybackslash}p{2.4cm}|}{List-II (About product mixture)} \\ \hline
(P) & \bfigi{82.2}{SA28_p2_chem_s2_q07_1.png} & (1) & Optically active \\ \hline
(Q) & \bfigi{79.7}{SA28_p2_chem_s2_q07_2.png} & (2) & Diastereoisomers \\ \hline
(R) & \bfigi{79.6}{SA28_p2_chem_s2_q07_3.png} & (3) & Optically inactive \\ \hline
(S) & \bfigi{83.3}{SA28_p2_chem_s2_q07_4.png} & (4) & Meso \\ \hline
\end{tabular}}\par
\begin{qoptions}
\opt{A}{P $\rightarrow$ 1,2;Q $\rightarrow$ 3;R $\rightarrow$ 3;S $\rightarrow$ 3,4}
\opt{B}{P $\rightarrow$ 1,3;Q $\rightarrow$ 2;R $\rightarrow$ 3;S $\rightarrow$ 3,4}
\opt{C}{P $\rightarrow$ 1,2;Q $\rightarrow$ 3;R $\rightarrow$ 1;S $\rightarrow$ 2,3}
\opt{D}{P $\rightarrow$ 1,2;Q $\rightarrow$ 4;R $\rightarrow$ 2;S $\rightarrow$ 3,4}
\end{qoptions}
\end{bankq}
\begin{bankq}[id=SA30-P1-CHEM-S1-Q05, source={Score Advanced - Test 30 (04 May 2024)}, paper=1, qno=5, type=single, subtopic={Optical}, answer={A}, answer_src=key, difficulty=4, diff_parts={K1 S2 I1 C1 F0}, diff_note={concepts: compound C in a reaction sequence | steps: follow the reagents from the starting material; identify intermediates; C structure}, also_in={Score Advanced - Test 31 (04 May 2024) P1 Q5}]
\bfigi{161.8}{SA30_p1_chem_s1_q05_R.png}\par
(C) is
\begin{qoptions}
\opt{A}{\bfigi{45.7}{SA30_p1_chem_s1_q05_A.png}}
\opt{B}{\bfigi{64.5}{SA30_p1_chem_s1_q05_B.png}}
\opt{C}{\bfigi{41.3}{SA30_p1_chem_s1_q05_C.png}}
\opt{D}{\bfigi{76.0}{SA30_p1_chem_s1_q05_D.png}}
\end{qoptions}
\end{bankq}
\begin{bankq}[id=SA32-P1-CHEM-S1-Q03, source={Score Advanced - Test 32 (08 May 2024)}, paper=1, qno=3, type=single, subtopic={Optical}, answer={A}, answer_src=key, difficulty=4, diff_parts={K1 S2 I1 C1 F0}, diff_note={concepts: most stable conformer of substituted molecules being achiral | steps: draw anti and gauche conformers; check symmetry of the lowest energy conformer for each compound; chirality}]
In which of the following cases, the most stable conformer is achiral?\par
\bfig{81.8}{SA32_p1_chem_s1_q03_1.png}
\begin{qoptions}
\opt{A}{Only IV}
\opt{B}{I and IV}
\opt{C}{I, II and III}
\opt{D}{In all the cases}
\end{qoptions}
\end{bankq}
\begin{bankq}[id=SA32-P1-CHEM-S3-Q03, source={Score Advanced - Test 32 (08 May 2024)}, paper=1, qno=3, type=numerical, subtopic={Optical}, answer={3}, answer_src=key, difficulty=5, diff_parts={K1 S2 I1 C1 F1}, diff_note={concepts: count correct stereochemical statements | steps: identify products X Y M N A B from the reactions; relation between each pair enantiomer diastereomer or identical; evaluate four statements}]
In the given reactions:-\par
\bfig{157.5}{SA32_p1_chem_s3_q03_1.png}
Number of correct statement(s) among the following are :-\par
(1) X and Y are enantiomers\par
(2) M and N are enantiomers\par
(3) A and B are diastereomers\par
(4) (X and Y) are identical with (A and B)\par
(5) (P and Q) are identical with (M and N)\par
(6) M and N are geometrical isomers
\end{bankq}
\begin{bankq}[id=SA33-P2-CHEM-S2-Q06, source={Score Advanced - Test 33 (12 May 2024)}, paper=2, qno=6, type=multi, subtopic={Optical}, answer={A,B,D}, answer_src=key, difficulty=5, diff_parts={K1 S2 I1 C1 F1}, diff_note={concepts: resolvable biphenyl or allene type compounds | steps: chirality requirement: ortho substituents big enough for hindered rotation in biphenyls; allene with two different groups on each end; check each drawn compound}, also_in={Score Advanced - Test 34 (12 May 2024) P2 Q6}]
Which of the following is/are resolvable ?
\begin{qoptions}
\opt{A}{\bfigi{55.0}{SA33_p2_chem_s2_q06_a.png}}
\opt{B}{\bfigi{67.6}{SA33_p2_chem_s2_q06_b.png}}
\opt{C}{\bfigi{36.7}{SA33_p2_chem_s2_q06_c.png}}
\opt{D}{\bfigi{53.0}{SA33_p2_chem_s2_q06_d.png}}
\end{qoptions}
\end{bankq}
\begin{bankq}[id=SA33-P2-CHEM-S3-Q08, source={Score Advanced - Test 33 (12 May 2024)}, paper=2, qno=8, type=numerical, subtopic={Optical}, answer={4}, answer_src=key, difficulty=5, diff_parts={K1 S2 I1 C1 F1}, diff_note={concepts: number of stereoisomers of the drawn compound | steps: count stereocentres and double bond geometries; check symmetry and meso forms; total from 2^n with corrections}, also_in={Score Advanced - Test 34 (12 May 2024) P2 Q8}]
How many stereoisomers are possible in \bfigi{36.3}{SA33_p2_chem_s3_q08_1.png}
\end{bankq}
\begin{bankq}[id=SA34-P1-CHEM-S3-Q01, source={Score Advanced - Test 34 (12 May 2024)}, paper=1, qno=1, type=single, subtopic={Structural}, answer={4}, answer_src=key, difficulty=5, diff_parts={K1 S2 I2 C1 F0}, diff_note={concepts: number of fractions from fractional distillation of monochloro derivatives | steps: enumerate distinct monochloro products of the given alkane; distinguish enantiomers (same boiling point) and diastereomers (different); count fractions | I: enantiomers cannot be separated by distillation}, also_in={Score Advanced - Test 36 (19 May 2024) P1 Q1}]
Total number of fractions obtained on fractional distillation of all monochloro derivatives product(s)\par
\bfig{38.5}{SA34_p1_chem_s3_q01_1.png}
\end{bankq}
\begin{bankq}[id=SA36-P2-CHEM-S4-Q10, source={Score Advanced - Test 36 (19 May 2024)}, paper=2, qno=10, type=numerical, subtopic={Geometrical}, answer={3}, answer_src=key, difficulty=5, diff_parts={K1 S2 I1 C1 F1}, diff_note={concepts: number of stereoisomers of compound X with the cis double bond fixed | steps: identify stereocentres and the other double bond; the fixed cis double bond removes its E Z contribution; count 2^n with symmetry corrections}, also_in={Score Advanced - Test 37 (19 May 2024) P2 Q10; TOAS - Test 05 (03 Apr 2022) P2 Q3}]
\bfigi{57.6}{SA36_p2_chem_s4_q09_1.png}\par
Identify total number of stereoisomer's if double bond is present in cis configuration in given compound (X)
\end{bankq}
\begin{bankpassage}[id=SA39-P2-CHEM-S2-G7-8, source={Score Advanced - Test 39 (20 Apr 2025)}, questions={7-8}]
1g of methane diffused in 20 seconds under certain conditions. Under the same conditions, $\sqrt{20}$g of an open chain hydrocarbon (A) diffused in 40 seconds. A 10 mg of sample of (A) consumes 8.40 ml of H$_2$ gas measured at 0$^\circ$C and 760 mm pressure.\par
One of the open chain isomers of (A) on ozonolysis gives only formaldehyde \& glyoxal.
\end{bankpassage}
\begin{bankq}[id=SA39-P2-CHEM-S2-Q07, source={Score Advanced - Test 39 (20 Apr 2025)}, paper=2, qno=7, type=single, subtopic={Structural}, answer={Bonus}, answer_src=solved, working={Printed key for this paper is incomplete and the figure/data needed to solve it was not recoverable; banked as Bonus.}, difficulty=4, diff_parts={K1 S2 I1 C1 F0}, diff_note={concepts: product that a six membered monocyclic isomer of A cannot give on reductive ozonolysis | steps: molecular formula of A; cyclohexene type isomers with substituents; ozonolysis cleaves the double bond giving dicarbonyls; check which listed product cannot form}, group=SA39-P2-CHEM-S2-G7-8]
A six membered stable monocyclic isomer of (A) on reductive ozonolysis can NOT give which of the following products
\begin{qoptions}
\opt{A}{OHC $-$ CHO}
\opt{B}{OHC$-$CH$_2-$CHO}
\opt{C}{OHC$-$CH$_2-$CH$_2-$CHO}
\opt{D}{OHC$-$CH$_2-$CH$_2-$CH$_2-$CHO}
\end{qoptions}
\end{bankq}
\begin{bankq}[id=SA39-P2-CHEM-S2-Q08, source={Score Advanced - Test 39 (20 Apr 2025)}, paper=2, qno=8, type=single, subtopic={Structural}, answer={Bonus}, answer_src=solved, working={Printed key for this paper is incomplete and the figure/data needed to solve it was not recoverable; banked as Bonus.}, difficulty=6, diff_parts={K1 S3 I1 C2 F0}, diff_note={concepts: incorrect statement about the isomers of A | steps: enumerate isomers of A (alkene or alkyne with given formula); geometrical isomerism; Tollens (terminal alkyne) white precipitate; optical activity; ozonolysis giving only acetone for the symmetric tetra substituted alkene; find the incorrect statement | C: scanning all isomers for each property}, group=SA39-P2-CHEM-S2-G7-8]
Identify the incorrect statement.
\begin{qoptions}
\opt{A}{One of the isomers of (A) exhibits geometrical isomerism}
\opt{B}{One of the isomers of (A) gives white precipitate with Tollen 's reagent}
\opt{C}{One of the isomers of (A) is optically active.}
\opt{D}{One of the isomers of (A) on reductive ozonolysis gives only acetone.}
\end{qoptions}
\end{bankq}
\begin{bankq}[id=SA40-P1-CHEM-S3-Q07, source={Score Advanced - Test 40 (27 Apr 2025)}, paper=1, qno=7, type=matching, subtopic={Structural}, answer={Bonus}, answer_src=solved, working={Figure/data needed was not recoverable and no key given; banked as Bonus.}, difficulty=5, diff_parts={K1 S2 I1 C1 F1}, diff_note={concepts: match isomer pairs with the type  isomerism | steps: positional chain functional and metamers definitions; classify each drawn pair; match}]
Correct matching from the two lists.\par
{\footnotesize\setlength{\tabcolsep}{1.5pt}\noindent\begin{tabular}{|c|L{4.2cm}|c|L{2.4cm}|}
\hline
 & \multicolumn{1}{>{\centering\arraybackslash}p{4.2cm}|}{List-I} & & \multicolumn{1}{>{\centering\arraybackslash}p{2.4cm}|}{List-II} \\ \hline
(P) & \bfigi{80.8}{SA40_p1_chem_s3_q07_1.png} & (1) & Positional isomers \\ \hline
(Q) & CH$_3$CH$_2$CH$_2$OH \& CH$_3$OCH$_2$CH$_3$ & (2) & Metamers \\ \hline
(R) & \bfigi{83.7}{SA40_p1_chem_s3_q07_2.png} & (3) & Chain isomers \\ \hline
(S) & \bfigi{79.4}{SA40_p1_chem_s3_q07_3.png} & (4) & Functional isomers \\ \hline
\end{tabular}}\par
\end{bankq}
\begin{bankq}[id=SA40-P1-CHEM-S3-Q08, source={Score Advanced - Test 40 (27 Apr 2025)}, paper=1, qno=8, type=matching, subtopic={Structural}, answer={P$\to$3; Q$\to$2; R$\to$4; S$\to$1}, answer_src=key, difficulty=6, diff_parts={K1 S3 I1 C1 F1}, diff_note={concepts: pi bonding types in NO3- dimethyl ether SO3  sulphate | steps: NO3- has p pi-p pi only; dimethyl ether has none; SO3 has both p pi-p pi and p pi-d pi in the usual description; sulphate only p pi-d pi; match}]
Correct matching from the two lists.\par
{\footnotesize\setlength{\tabcolsep}{1.5pt}\noindent\begin{tabular}{|c|L{1.6cm}|c|L{4.8cm}|}
\hline
 & \multicolumn{1}{>{\centering\arraybackslash}p{1.6cm}|}{List-I} & & \multicolumn{1}{>{\centering\arraybackslash}p{4.8cm}|}{List-II (Type of $\pi-$bonds)} \\ \hline
(P) & NO$_3^-$ & (1) & Only p$\pi-$d$\pi$ bonding \\ \hline
(Q) & (CH$_3$)$_2$O & (2) & No p$\pi-$d$\pi$ and p$\pi-$p$\pi$ bonding \\ \hline
(R) & SO$_3$ & (3) & Only p$\pi-$p$\pi$ bonding \\ \hline
(S) & SO$_4^{2-}$ & (4) & Both p$\pi-$p$\pi$ and p$\pi-$d$\pi$ bonding \\ \hline
\end{tabular}}\par
\end{bankq}
\begin{bankq}[id=SA40-P1-CHEM-S4-Q02, source={Score Advanced - Test 40 (27 Apr 2025)}, paper=1, qno=2, type=numerical, subtopic={Geometrical}, answer={5}, answer_src=key, difficulty=5, diff_parts={K1 S2 I1 C1 F1}, diff_note={concepts: number of alkenes exhibiting geometrical isomerism | steps: geometric isomerism needs different substituents on both ends; check each drawn alkene; (CH3)2C=CHMe has two identical groups on one end so none; count}]
How many of these can exhibit geometrical isomerism?\par
\bfigi{59.0}{SA40_p1_chem_s4_q02_2.png}\quad \bfigi{50.8}{SA40_p1_chem_s4_q02_3.png}\par \bfigi{63.0}{SA40_p1_chem_s4_q02_1.png}\quad \bfigi{52.9}{SA40_p1_chem_s4_q02_4.png}\quad \bfigi{64.4}{SA40_p1_chem_s4_q02_5.png}\par (vi) (CH$_3$)$_2$C=CHMe\quad \bfigi{65.8}{SA40_p1_chem_s4_q02_7.png}\par \bfigi{65.1}{SA40_p1_chem_s4_q02_8.png}\quad \bfigi{63.0}{SA40_p1_chem_s4_q02_6.png}
\end{bankq}

\section*{SRGBT papers}
\begin{bankq}[id=SRGBT07-P1-CHEM-S1-Q23, source={SRGBT - Test 07 (05 Sep 2021)}, paper=1, qno=23, type=multi, subtopic={Structural}, answer={D}, answer_src=key, difficulty=3, diff_parts={K1 S1 I1 C0 F1}, diff_note={concepts: conformational stability with intramolecular hydrogen bonding and gauche interactions | steps: the gauche form allowing the H-bond is the most stable so D}]
The most stable conformation of given compound.\par
\bfig{83.6}{SRGBT07_p1_chem_s0_q23_1.png}
\begin{qoptions}
\opt{A}{\bfigi{111.7}{SRGBT07_p1_chem_s0_q23_3.png}}
\opt{B}{\bfigi{93.2}{SRGBT07_p1_chem_s0_q23_2.png}}
\opt{C}{\bfigi{111.5}{SRGBT07_p1_chem_s0_q23_5.png}}
\opt{D}{\bfigi{88.2}{SRGBT07_p1_chem_s0_q23_4.png}}
\end{qoptions}
\end{bankq}
\begin{bankq}[id=SRGBT07-P1-CHEM-S1-Q37, source={SRGBT - Test 07 (05 Sep 2021)}, paper=1, qno=37, type=multi, subtopic={Structural}, answer={6}, answer_src=key, difficulty=6, diff_parts={K1 S3 I1 C1 F1}, diff_note={concepts: enumerating saturated bicyclic skeletons of C8H14 | steps: no primary H means no methyl groups so all carbons are in rings; no quaternary carbon limits bridgeheads to CH; list bicyclo[x.y.z] with x + y + z = 6; include fused and bridged types; count 6}]
How many saturated bicyclo hydrocarbon are possible having molecular formula C$_8$H$_{14}$ with no primary hydrogen and no quaternary carbon?
\end{bankq}
\begin{bankq}[id=SRGBT14-P1-CHEM-S1-Q31, source={SRGBT - Test 14 (24 Oct 2021)}, paper=1, qno=31, type=multi, subtopic={Structural}, answer={A,C}, answer_src=key, difficulty=2, diff_parts={K1 S1 I0 C0 F1}, diff_note={concepts: conformers differ only by rotation about single bonds | steps: A and C are the same compound rotated about a single bond}]
Which of the following is/are conformational of (p) ?\par
\bfigi{83.6}{SRGBT14_p1_chem_s0_q31_1.png}
\begin{qoptions}
\opt{A}{\bfigi{61.8}{SRGBT14_p1_chem_s0_q31_A.png}}
\opt{B}{\bfigi{75.0}{SRGBT14_p1_chem_s0_q31_B.png}}
\opt{C}{\bfigi{74.1}{SRGBT14_p1_chem_s0_q31_C.png}}
\opt{D}{\bfigi{64.6}{SRGBT14_p1_chem_s0_q31_D.png}}
\end{qoptions}
\end{bankq}
\begin{bankq}[id=SRGBT14-P1-CHEM-S1-Q32, source={SRGBT - Test 14 (24 Oct 2021)}, paper=1, qno=32, type=multi, subtopic={Structural}, answer={B,D}, answer_src=key, difficulty=2, diff_parts={K1 S1 I0 C0 F1}, diff_note={concepts: positional chain and conformational isomer relations | steps: compare connectivity in each pair; B is not identical and D is not a chain isomer}]
Select incorrect relations.
\begin{qoptions}
\opt{A}{\bfigi{99.4}{SRGBT14_p1_chem_s0_q32_A.png} $\Rightarrow$ Positional isomer}
\opt{B}{\bfigi{119.5}{SRGBT14_p1_chem_s0_q32_B.png} $\Rightarrow$ Identical}
\opt{C}{\bfigi{85.0}{SRGBT14_p1_chem_s0_q32_C.png} $\Rightarrow$ Conformational isomers}
\opt{D}{\bfigi{117.4}{SRGBT14_p1_chem_s0_q32_D.png} $\Rightarrow$ Chain isomer}
\end{qoptions}
\end{bankq}
\begin{bankpassage}[id=SRGBT14-P1-CHEM-S1-G33-34, source={SRGBT - Test 14 (24 Oct 2021)}, questions={33-34}]
Consider the following molecules having same molecular formula.\par
\bfig{302.0}{SRGBT14_p1_chem_s0_q33_1.png}
\end{bankpassage}
\begin{bankq}[id=SRGBT14-P1-CHEM-S1-Q33, source={SRGBT - Test 14 (24 Oct 2021)}, paper=1, qno=33, type=multi, subtopic={Structural}, answer={B}, answer_src=key, difficulty=2, diff_parts={K1 S1 I0 C0 F1}, diff_note={concepts: types of structural isomerism | steps: the compounds given differ in chain and functional group and metamerism but not in position}, group=SRGBT14-P1-CHEM-S1-G33-34]
Which isomerism is not observed in between above compounds.
\begin{qoptions}
\opt{A}{Chain Isomer}
\opt{B}{Position Isomer}
\opt{C}{Metamer Isomer}
\opt{D}{Functional Isomer}
\end{qoptions}
\end{bankq}
\begin{bankq}[id=SRGBT14-P1-CHEM-S1-Q34, source={SRGBT - Test 14 (24 Oct 2021)}, paper=1, qno=34, type=multi, subtopic={Structural}, answer={B}, answer_src=key, difficulty=2, diff_parts={K1 S1 I0 C0 F1}, diff_note={concepts: acidity of the listed molecules | steps: c has the most stabilised conjugate base}, group=SRGBT14-P1-CHEM-S1-G33-34]
Most acidic in the above molecules.
\begin{qoptions}
\opt{A}{d}
\opt{B}{c}
\opt{C}{f}
\opt{D}{e}
\end{qoptions}
\end{bankq}
\begin{bankq}[id=SRGBT14-P1-CHEM-S1-Q39, source={SRGBT - Test 14 (24 Oct 2021)}, paper=1, qno=39, type=multi, subtopic={Structural}, answer={3}, answer_src=key, difficulty=2, diff_parts={K1 S1 I0 C0 F1}, diff_note={concepts: conformations of butane about C2-C3 | steps: distinct energy minima anti and two gauche so 3}]
How many conformers are possible for butane. (Consider rotation about C$_2$--C$_3$ bond)
\end{bankq}

\section*{Nurture SRG papers}
\begin{bankq}[id=NSRG03-P1-CHEM-S1-Q02, source={Nurture SRG - Test 03 (04 Feb 2024)}, paper=1, qno=2, type=single, subtopic={Optical}, answer={D}, answer_src=key, difficulty=1, diff_parts={K1 S1 I0 C0 F0}, diff_note={concepts: matching of reactions with named processes | steps: each of four items matched}]
Match list I with list II and select the correct answer from the given below the lists :\par
{\footnotesize\setlength{\tabcolsep}{2pt}\begin{tabular}{|c|c|c|l|}
\hline
\multicolumn{2}{|c|}{List - I} & & List - II\\
\hline
P. & \bfigi{78.6}{NSRG03_p1_chem_s1_q02_1.png} & i. & $(2R,3E)-3-$penten$-2-$ol\\
\hline
Q. & \bfigi{78.2}{NSRG03_p1_chem_s1_q02_2.png} & ii. & $(2R,3Z)-3-$penten$-2-$ol\\
\hline
R. & \bfigi{78.1}{NSRG03_p1_chem_s1_q02_3.png} & iii. & $(2S,3E)-3-$penten$-2-$ol\\
\hline
S. & \bfigi{79.8}{NSRG03_p1_chem_s1_q02_4.png} & iv. & $(2S,3Z)-3-$penten$-2-$ol\\
\hline
\end{tabular}}
\begin{qoptions}
\opt{A}{P--i, Q--iii, R--iv, S--ii}
\opt{B}{P--i, Q--iii, R--ii, S--iv}
\opt{C}{P--iv, Q--iii, R--i, S--ii}
\opt{D}{P--iv, Q--iii, R--ii, S--i}
\end{qoptions}
\end{bankq}
\begin{bankq}[id=NSRG03-P1-CHEM-S2-Q11, source={Nurture SRG - Test 03 (04 Feb 2024)}, paper=1, qno=11, type=multi, subtopic={Optical}, answer={A,C,D}, answer_src=key, difficulty=6, diff_parts={K2 S2 I1 C1 F1}, diff_note={concepts: multi step reaction, stereoisomers | steps: identify Q; identify R; stereoisomers of R with compensation; check each | I: internal and external compensation}]
Which of the following is/are correct statements ?\par
\bfig{180.7}{NSRG03_p1_chem_s2_q11_1.png}
\begin{qoptions}
\opt{A}{$Q$ is \bfigi{78.7}{NSRG03_p1_chem_s2_q11_2.png}}
\opt{B}{Stereomers of product $R$ involves internal as well as external compensation of optical rotation}
\opt{C}{Single enantiomerically pure form of $Q$ is formed}
\opt{D}{Enantiomers as well as diastereomers are formed as $R$}
\end{qoptions}
\end{bankq}
\begin{bankq}[id=NSRG03-P1-CHEM-S2-Q12, source={Nurture SRG - Test 03 (04 Feb 2024)}, paper=1, qno=12, type=multi, subtopic={Optical}, answer={B,C}, answer_src=key, difficulty=6, diff_parts={K2 S2 I1 C1 F1}, diff_note={concepts: multi step organic problem, stereocentres | steps: identify A B C from the figure; count stereogenic centres; compare | I: stereocentres created in each step}]
\bfig{230.6}{NSRG03_p1_chem_s2_q12_1.png}
\begin{qoptions}
\opt{A}{$A$ is \bfigi{96.7}{NSRG03_p1_chem_s2_q12_2.png}}
\opt{B}{$B$ is \bfigi{43.5}{NSRG03_p1_chem_s2_q12_3.png}}
\opt{C}{$C$ has 3 stereogenic centres}
\opt{D}{$B$ has the same number of stereogenic centres as $C$}
\end{qoptions}
\end{bankq}
\begin{bankq}[id=NSRG03-P2-CHEM-S3-Q03, source={Nurture SRG - Test 03 (04 Feb 2024)}, paper=2, qno=3, type=numerical, subtopic={Geometrical}, answer={10}, answer_src=solved, working={$x=2$ (geometrical isomers), $y=4$, $z=4$}, difficulty=5, diff_parts={K1 S2 I1 C1 F1}, diff_note={concepts: isomer counting | steps: diastereomers of the alkene; isomeric alcohols C4H10O; third part; combine | I: count each independent set}]
Number of diastereomers of $\mathrm{C_2(F)(C\ell)(Br)(I)}$ is $\mathbf x$,\par
Number of isomeric alcohols of $\mathrm{C_4H_{10}O}$ is $\mathbf y$ and\par
Number of isomeric esters of $\mathrm{C_4H_8O_2}$ is $\mathbf z$.\par
Calculate the value of $(\mathbf x+\mathbf y+\mathbf z)$.
\end{bankq}

\section*{Pre-Score (2023-24) papers}
\begin{bankq}[id=PS03-P2-CHEM-S3-Q02, source={Pre-Score (2023-24) - Test 03 (22 Oct 2023)}, paper=2, qno=2, type=numerical, subtopic={Optical}, answer={6.00}, answer_src=key, working={Printed key.}, difficulty=7, diff_parts={K1 S2 I2 C2 F1}, diff_note={concepts: isomers of C8H10O2 with phenol and optical activity | steps: degrees of unsaturation 4; phenolic OH with a chiral side chain; enumerate structures with FeCl3 colour and optical activity including enantiomers; count | I: combine phenol requirement with a stereocentre | C: enumerating isomers and stereoisomers}]
An organic compound (C$_8$H$_{10}$O$_2$) rotates plane-polarized light. It produces voilet color with neutral FeCl$_3$ solution. What is the total number of all the possible isomers (including stereoisomers) for this compound?
\end{bankq}
\begin{bankq}[id=PS08-P1-CHEM-S3-Q11, source={Pre-Score (2023-24) - Test 08 (03 Dec 2023)}, paper=1, qno=11, type=single, subtopic={Optical}, answer={A}, answer_src=key, working={Printed key.}, difficulty=5, diff_parts={K1 S2 I2 C1 F0}, diff_note={concepts: stereochemical outcome of four reactions | steps: identify mechanism and stereochemistry for each; count products and optical activity; match with list II | I: stereochemistry of addition or substitution | C: four reactions}]
\par\nobreak
\begin{center}\begin{tabular}{|L{0.8cm}|L{5cm}|L{0.8cm}|L{4cm}|}\hline
\multicolumn{2}{|c|}{\begin{tabular}[c]{@{}c@{}}\textbf{List-I}\\\textbf{(For the given reaction)}\end{tabular}} & \multicolumn{2}{c|}{\begin{tabular}[c]{@{}c@{}}\textbf{List-II}\\\textbf{(For final product or product mixture)}\end{tabular}}\\\hline
(P) & \bfigi{77.9}{PS08_p1_chem_s3_q10_1.png} & (1) & Optically active product mixture\\\hline
(Q) & \bfigi{112.3}{PS08_p1_chem_s3_q10_2.png} & (2) & Optically inactive product / product mixture\\\hline
(R) & \bfigi{108.0}{PS08_p1_chem_s3_q10_3.png} & (3) & Product can exhibit optical isomerism\\\hline
(S) & \bfigi{94.6}{PS08_p1_chem_s3_q10_4.png} & (4) & Carbocation is involved in the reaction sequence\\\hline
 & & (5) & Racemic mixture is obtained\\\hline
\end{tabular}\end{center}
\begin{qoptions}
\opt{A}{P $\to$ 2,4; Q $\to$ 1,3; R $\to$ 2,3,4; S $\to$ 2,3,4,5}
\opt{B}{P $\to$ 1,2,3; Q $\to$ 1,3,4; R $\to$ 1,2; S $\to$ 4,5}
\opt{C}{P $\to$ 4,5; Q $\to$ 1,2,3; R $\to$ 3,4; S $\to$ 2,4}
\opt{D}{P $\to$ 1,2,3; Q $\to$ 1,4; R $\to$ 2,5; S $\to$ 1,2,3}
\end{qoptions}
\end{bankq}
\begin{bankq}[id=PS08-P2-CHEM-S4-Q05, source={Pre-Score (2023-24) - Test 08 (03 Dec 2023)}, paper=2, qno=5, type=numerical, subtopic={Structural}, answer={5}, answer_src=key, working={Printed key.}, difficulty=5, diff_parts={K1 S1 I2 C1 F1}, diff_note={concepts: cyclic structural isomers of C5H10 | steps: count cycloalkanes with five carbons: cyclopentane, methylcyclobutane, 1,1 and 1,2 and 1,3 dimethylcyclopropane, ethylcyclopropane | I: enumerate rings and substituents systematically | C: counting isomers}]
The total number of cyclic, structural isomers for a compound with molecular formula C$_5$H$_{10}$ is \underline{\hspace{1cm}}
\end{bankq}

\section*{Enthusiast Pre-Score papers}
\begin{bankq}[id=EPS02-P2-CHEM-S3-Q01, source={Enthusiast Pre-Score - Test 02 (11 Nov 2020)}, paper=2, qno=1, type=numerical, subtopic={Structural}, answer={5.00}, answer_src=key, difficulty=4, diff_parts={K1 S2 I1 C0 F1}, diff_note={concepts: reduction of carboxylic acid isomers | steps: list isomeric acids; those giving primary alcohol; count}]
Identify total number of carboxylic acids with C$_5$H$_{10}$O$_2$ molecular formula gives 1$^\circ$-Alcohol on reaction with LiAlH$_4$.
\end{bankq}

\section*{JEE Advanced (2021 Series) papers}
\begin{bankq}[id=JA02-P1-CHEM-S1-Q04, source={JEE Advanced (2021 Series) - Test 02 (04 Feb 2021)}, paper=1, qno=4, type=single, subtopic={Structural}, answer={C}, answer_src=key, difficulty=4, diff_parts={K1 S2 I1 C1 F0}, diff_note={concepts: identical molecules from Newman projections | steps: convert each Newman projection to a stereo description by assigning R or S at the front and back carbons; rotate about the C-C bond to compare; identify which pairs are the same molecule}]
Newman projections \textbf{P}, \textbf{Q}, \textbf{R} and \textbf{S} are shown below :
\bfig{395.5}{JA02_p1_chem_s1_q04_1.png}

Which one of the following options represents identical molecules ?
\begin{qoptions}
\opt{A}{P and Q}
\opt{B}{Q and S}
\opt{C}{Q and R}
\opt{D}{R and S}
\end{qoptions}
\end{bankq}
\begin{bankq}[id=JA02-P1-CHEM-S1-Q06, source={JEE Advanced (2021 Series) - Test 02 (04 Feb 2021)}, paper=1, qno=6, type=single, subtopic={Optical}, answer={C}, answer_src=key, difficulty=5, diff_parts={K1 S3 I1 C1 F0}, diff_note={concepts: relationship between Fischer projections of sugar stereoisomers | steps: assign configuration at both stereocentres of D-erythrose; do the same for P Q R S by swapping groups and counting exchanges; same at both centres is identical; opposite at both is enantiomer; one centre different is diastereomer; match codes}]
The Fischer projection of D-erythrose is shown below.
\begin{center}\bfigi{55.5}{JA02_p1_chem_s1_q06_1.png}\\ D-Erythrose\end{center}
D-Erythrose and its isomers are listed as \textbf{P}, \textbf{Q}, \textbf{R} and \textbf{S} in \textbf{Column-I}. Choose the correct relationship of \textbf{P}, \textbf{Q}, \textbf{R} and \textbf{S} with D-erythrose from \textbf{Column-II}.\par\nobreak
\begin{center}
\begin{tabular}{@{}l@{\hspace{4pt}}c@{\hspace{4pt}}l@{}}
\textbf{Column-I} & & \textbf{Column-II}\\[2pt]
\textbf{(P)} & \bfigi{109.3}{JA02_p1_chem_s1_q06_2.png} & 1.\ Diastereomer\\[4pt]
\textbf{(Q)} & \bfigi{106.0}{JA02_p1_chem_s1_q06_3.png} & 2.\ Identical\\[4pt]
\textbf{(R)} & \bfigi{109.3}{JA02_p1_chem_s1_q06_4.png} & 3.\ Enantiomer\\[4pt]
\textbf{(S)} & \bfigi{106.0}{JA02_p1_chem_s1_q06_5.png} &
\end{tabular}
\end{center}
\begin{qoptions}
\opt{A}{P$\to$2, Q$\to$3, R$\to$2, S$\to$2}
\opt{B}{P$\to$3, Q$\to$1, R$\to$1, S$\to$2}
\opt{C}{P$\to$2, Q$\to$1, R$\to$1, S$\to$3}
\opt{D}{P$\to$2, Q$\to$3, R$\to$3, S$\to$1}
\end{qoptions}
\end{bankq}
\begin{bankq}[id=JA03-P2-CHEM-S1-Q03, source={JEE Advanced (2021 Series) - Test 03 (01 Mar 2021)}, paper=2, qno=3, type=multi, subtopic={Structural}, answer={A,B,C,D}, answer_src=key, difficulty=6, diff_parts={K1 S3 I1 C1 F1}, diff_note={concepts: counting open chain and cyclic isomers of C5H6 with stereochemistry | steps: degree of unsaturation 3; enumerate open chain conjugated alkadienes; list cyclic isomers and mark chiral centres; count open chain geometrical stereoisomers; count optically active ones; check each statement | C: enumeration of many isomers with stereo variants}]
True/correct statement for an organic compound having Molecular formula C$_5$H$_6$ is/are
\begin{qoptions}
\opt{A}{Total open chain conjugated alkadiene possible is $2$}
\opt{B}{Total possible cyclic isomer with chiral centre is $2$}
\opt{C}{Total open chain stereoisomers possible is $4$}
\opt{D}{Total open chain optically active isomers possible is $2$}
\end{qoptions}
\end{bankq}
\begin{bankq}[id=JA04-P2-CHEM-S1-Q03, source={JEE Advanced (2021 Series) - Test 04 (04 Apr 2021)}, paper=2, qno=3, type=multi, subtopic={Optical}, answer={A,B,D}, answer_src=key, difficulty=4, diff_parts={K1 S2 I1 C0 F1}, diff_note={concepts: resolution of racemic carboxylic acids | steps: only chiral acids form diastereomeric salts with an enantiopure amine; identify the chiral centres in each drawn acid; pick those that are chiral}]
Which of the following carboxylic acid could be resolved by reaction with an enantiometrically pure chiral amine ?
\begin{qoptions}
\opt{A}{\bfigi{77.3}{JA04_p2_chem_s1_q03_1.png}}
\opt{B}{\bfigi{88.5}{JA04_p2_chem_s1_q03_B.png}}
\opt{C}{\bfigi{123.1}{JA04_p2_chem_s1_q03_C.png}}
\opt{D}{\bfigi{73.7}{JA04_p2_chem_s1_q03_D.png}}
\end{qoptions}
\end{bankq}
\begin{bankq}[id=JA06-P1-CHEM-S3-Q07, source={JEE Advanced (2021 Series) - Test 06 (22 Apr 2021)}, paper=1, qno=7, type=single, subtopic={Structural}, answer={1.60}, answer_src=key, difficulty=7, diff_parts={K1 S3 I2 C2 F0}, diff_note={concepts: positional optical and geometrical isomerism of alkynes | steps: enumerate alkynes of C3FClBrI with the triple bond at terminal or internal positions; those with a stereocentre give enantiomers; count p q r s separately; fractional distillation counts distinct physical species; evaluate the expression | I: classify each isomer by chirality before counting distillation fractions | C: lengthy enumeration of isomers with stereo variants}]
Various isomeric alkynes have molecular formula $\mathrm{C_3FClBrI}$. For such alkynes if\\
$p=$ Number of positional isomers\\
$q=$ Number of optically active isomers\\
$r=$ Number of fractions obtained on fractional distillation of a mixture of all possible alkynes\\
$s=$ Number of diastereomeric pairs\\
Then find the value of $\left(\dfrac{p+q-r+s}{5}\right)$.
\end{bankq}
\begin{bankq}[id=JA07-P2-CHEM-S3-Q06, source={JEE Advanced (2021 Series) - Test 07 (25 Apr 2021)}, paper=2, qno=6, type=numerical, subtopic={Optical}, answer={3.00}, answer_src=key, difficulty=6, diff_parts={K1 S2 I2 C1 F1}, diff_note={concepts: optically active stereoisomers of a molecule with several stereocentres | steps: count stereocentres and any symmetry in the drawn molecule; total stereoisomers 2^n adjusted for meso forms; subtract the meso or achiral ones to get optically active isomers = 3 | I: symmetric ends reduce the isomer count through meso forms}]
Calculate total number of optically active stereoisomers for following compound?
\bfig{177.6}{JA07_p2_chem_s3_q06_1.png}
\end{bankq}
\begin{bankq}[id=JA08-P1-CHEM-S1-Q02, source={JEE Advanced (2021 Series) - Test 08 (01 May 2021)}, paper=1, qno=2, type=multi, subtopic={Optical}, answer={B,C,D}, answer_src=solved, working={(A) none of the C5H10O ketones is chiral; (B) 3-methylhexan-3-ol gives 5 alkenes (1+2+2 with E/Z); (C) SN1 of PhCH(Cl)Et gives two enantiomeric ethers; (D) Markovnikov HBr gives a product with 3 stereocentres (D, T, new C) = 8.}, difficulty=5, diff_parts={K1 S2 I1 C1 F1}, diff_note={concepts: counting products from stereoisomeric and regioisomeric reactions | steps: identify chiral ketones of C5H10O; dehydration of 3-methylhexan-3-ol with E/Z alkenes; SN1 ether products as enantiomers; HBr addition to an alkene with D and T labels giving three stereocentres}]
Identify correct match for total number of products ?
\begin{qoptions}
\opt{A}{Optically active ketones with molecular formula [C$_5$H$_{10}$O] $\xrightarrow{\text{LiAlH}_4}$ Total six alcohol products are formed}
\opt{B}{\bfigi{52.0}{JA08_p1_chem_s1_q02_B.png} $\xrightarrow[\Delta]{\text{Conc. H}_2\text{SO}_4}$ Five alkenes are formed during reaction}
\opt{C}{\bfigi{45.5}{JA08_p1_chem_s1_q02_C.png} $\xrightarrow[(\mathrm{S_N}1)]{\text{EtOH}}$ Two ether products are formed}
\opt{D}{\bfigi{51.4}{JA08_p1_chem_s1_q02_D.png} $\xrightarrow[\text{CCl}_4]{\text{HBr}}$ Total eight mono bromo major products are formed}
\end{qoptions}
\end{bankq}
\begin{bankq}[id=JA08-P1-CHEM-S3-Q01, source={JEE Advanced (2021 Series) - Test 08 (01 May 2021)}, paper=1, qno=1, type=numerical, subtopic={Geometrical}, answer={5.00}, answer_src=solved, working={two C=C each E/Z, two CH(OH)CH3 centres R/S and the central carbon: counting arrangements of ligands E,Z,R,S gives $1+1+1+1+1+1+1+1+2=10$ stereoisomers, so $X-5=5$.}, difficulty=7, diff_parts={K1 S3 I1 C2 F1}, diff_note={concepts: stereoisomers from double bonds and chiral centres | steps: two C=C each give E/Z; two CH(OH)CH3 centres give R/S; central carbon with two identical types becomes a pseudo asymmetric centre; count arrangements of ligands E Z R S at the central carbon; total 10 | C: pseudoasymmetric centre counting with enantiomorphic ligands}]
Total number of stereoisomers for following structure is $X$. Then identify the value of $(X-5)$ is :-
\bfig{121.8}{JA08_p1_chem_s3_q01_1.png}
\end{bankq}
\begin{bankq}[id=JA09-P1-CHEM-S1-Q03, source={JEE Advanced (2021 Series) - Test 09 (07 May 2021)}, paper=1, qno=3, type=single, subtopic={Optical}, answer={B}, answer_src=key, difficulty=6, diff_parts={K2 S3 I1 C1 F0}, diff_note={concepts: Gabriel synthesis stereochemistry, amine formation from nitrogen compounds by reduction | steps: check reaction (i) which is a Grignard reaction and gives no amine; reaction (ii) is a quinone imine reduction giving an achiral product; reaction (iii) Gabriel reagent with the chiral bromide gives one inverted chiral amine after hydrolysis; reaction (iv) isocyanide hydrolysis gives a secondary amine not a primary one; reaction (v) Na with ethanol gives a racemic or achiral amine; only (iii) qualifies | C: five separate reactions each judged for amine and optical activity}]
Identify which reaction(s) produce(s) single optically active Amine molecule as product\\
\bfigi{88.2}{JA09_p1_chem_s1_q03_1.png} $\xrightarrow[\text{(1eq.)}]{\text{Mg/Ether}}$ \quad Optically pure\\
\bfigi{48.9}{JA09_p1_chem_s1_q03_2.png} $\xrightarrow[\text{(ii) H}_2/\text{Ni (excess)}/\Delta]{\text{(i) conc.HBr}}$\\
\bfigi{91.8}{JA09_p1_chem_s1_q03_3.png} $+$ \bfigi{72.2}{JA09_p1_chem_s1_q03_4.png}\\
\bfigi{79.8}{JA09_p1_chem_s1_q03_5.png} $\xrightarrow{\text{Na/Et--OH}}$
\begin{qoptions}
\opt{A}{ii, iii, iv}
\opt{B}{iii only}
\opt{C}{iii , iv}
\opt{D}{none of these}
\end{qoptions}
\end{bankq}
\begin{bankq}[id=JA09-P2-CHEM-S1-Q01, source={JEE Advanced (2021 Series) - Test 09 (07 May 2021)}, paper=2, qno=1, type=multi, subtopic={Optical}, answer={A,B,D}, answer_src=key, difficulty=5, diff_parts={K1 S3 I1 C0 F1}, diff_note={concepts: stereo relationship between pairs of optical isomers | steps: assign R or S to each stereocentre in I II III and IV; compare I with II and find mirror images; compare III with IV and find mirror images; compare I with IV and find non mirror images with one inverted centre; judge each option from these relations | I: use the pattern of inverted centres to separate enantiomers from diastereomers}]
\bfig{344.9}{JA09_p2_chem_s1_q01_1.png}

Which of the following statement are true about these isomers ?
\begin{qoptions}
\opt{A}{I and II are a pair of enantiomer}
\opt{B}{III and IV are a pair of enantiomer}
\opt{C}{I and II are a pair of diastereomers}
\opt{D}{I and IV are a pair of diastereomers}
\end{qoptions}
\end{bankq}
\begin{bankq}[id=JA11-P1-CHEM-S3-Q04, source={JEE Advanced (2021 Series) - Test 11 (20 May 2021)}, paper=1, qno=4, type=single, subtopic={Optical}, answer={8}, answer_src=key, difficulty=5, diff_parts={K2 S2 I1 C1 F0}, diff_note={concepts: axial chirality of allenes and cumulenes, geometrical isomerism on a three membered ring | steps: identify the stereogenic units in the drawn structure; count the allene axes and the ring positions with different substituents; combine independent units as powers of two while removing duplicates from symmetry; result is 8 | I: each allene or cumulene unit gives its own chirality element beyond the ring centres | C: several kinds of stereo elements counted together}, also_in={JEE Advanced (2021 Series) - Test 15 (12 Jun 2021) P1 Q4; Nurture (General) - Test 01 (12 Jun 2021) P1 Q4}]
Total number of stereo isomers corresponding to structure
\bfig{128.1}{JA11_p1_chem_s3_q04_1.png}
\end{bankq}
\begin{bankq}[id=JA12-P1-CHEM-S2-Q02, source={JEE Advanced (2021 Series) - Test 12 (25 May 2021)}, paper=1, qno=2, type=numerical, subtopic={Optical}, answer={6}, answer_src=solved, working={Two equal stereocentres give three stereoisomers (RR, SS, meso) for each of E and Z alkene: 6.}, difficulty=4, diff_parts={K1 S2 I1 C0 F1}, diff_note={concepts: stereoisomer count with symmetry of a symmetrical diene | steps: there are two identical stereocentres and one double bond; for each of E and Z the configurations RR, SS and meso occur; total six | I: identical ends make the RS form meso so 8 reduces to 6}]
\bfig{100.7}{JA12_p1_chem_s2_q02_1.png}
Give total number of stereoisomers.
\end{bankq}
\begin{bankq}[id=JA13-P1-CHEM-S1-Q02, source={JEE Advanced (2021 Series) - Test 13 (31 May 2021)}, paper=1, qno=2, type=multi, subtopic={Optical}, answer={B}, answer_src=solved, working={Case I (ethene) gives only E/Z isomers and no enantiomers; case II gives enantiomers and diastereomers so A is wrong; mixing a dextro and a laevo molecule of different diastereomers does not necessarily cancel; total 6 + 40 = 46 not 86.}, difficulty=6, diff_parts={K1 S3 I1 C1 F1}, diff_note={concepts: counting stereoisomers of fully differently substituted ethene and ethane | steps: ethene with four different groups gives only geometrical isomers and no enantiomers; ethane with six different groups has two stereocentres giving enantiomers and diastereomers; count ten ways of choosing three groups for the first carbon with four stereoisomers each giving 40; total with ethene (6) is 46; so only the statement about enantiomers holds | I: identical constitutions from swapping the two carbons must be counted once | C: counting the arrangements of six different groups}]
Carefully study the following 2 cases and choose the correct statements followed ?\\
Case \# I. All H atoms in ethene are substituted differently without significantly distorting the original geoemetry.\\
Case \# II. All H atoms in ethane are substituted differently without significantly distorting the original geoemetry.\\
Choose the correct statements(s) from the following :
\begin{qoptions}
\opt{A}{Diastereomers will be produced only in Case I , not in Case II.}
\opt{B}{Enantiomers will be produced only in Case II, not in Case I}
\opt{C}{Mixing equal weights of one dextro and one laevo molecule under Case II will necessarily yield an optically inactive mixture}
\opt{D}{Total number of space isomers possible under both cases $=86$}
\end{qoptions}
\end{bankq}
\begin{bankq}[id=JA14-P2-CHEM-S3-Q10, source={JEE Advanced (2021 Series) - Test 14 (06 Jun 2021)}, paper=2, qno=10, type=single, subtopic={Geometrical}, answer={4}, answer_src=key, working={Printed key.}, difficulty=6, diff_parts={K1 S2 I2 C2 F0}, diff_note={concepts: stereoisomers of hydrazones and oximes | steps: identify the products with C=N-NH from the reagent; each C=N gives E and Z; add chiral centres; count total stereoisomers | I: each C=N contributes geometric isomerism plus any chiral centres | C: counting stereo isomers across multiple products}]
The number of stereoisomers of organic products containing $-$C=N--NH$-$ obtained in the following reaction in acidic medium.
\bfig{125.6}{JA14_p2_chem_s3_q10_1.png}
\end{bankq}

\section*{TOAS papers}
\begin{bankq}[id=TOAS01-P1-CHEM-S1-Q03, source={TOAS - Test 01 (30 Jan 2022)}, paper=1, qno=3, type=single, subtopic={Optical}, answer={D}, answer_src=key, difficulty=4, diff_parts={K2 S2 I1 C0 F0}, diff_note={concepts: chirality, symmetry elements | steps: test the plane of symmetry of each structure; find the one without any symmetry element; check the chirality centre | I: absence of any improper axis}]
Which of following is optically active in nature
\begin{qoptions}
\opt{A}{\bfigi{89.6}{TOAS01_p1_chem_s1_q03_A.png}}
\opt{B}{\bfigi{102.3}{TOAS01_p1_chem_s1_q03_B.png}}
\opt{C}{\bfigi{44.1}{TOAS01_p1_chem_s1_q03_C.png}}
\opt{D}{\bfigi{51.7}{TOAS01_p1_chem_s1_q03_D.png}}
\end{qoptions}
\end{bankq}
\begin{bankq}[id=TOAS01-P1-CHEM-S1-Q04, source={TOAS - Test 01 (30 Jan 2022)}, paper=1, qno=4, type=single, subtopic={Optical}, answer={B}, answer_src=key, difficulty=3, diff_parts={K2 S1 I1 C0 F0}, diff_note={concepts: isomerism, structure comparison | steps: classify each pair; identify the pair that is truly isomeric}]
Which one is correct in a given isomeric pair :
\begin{qoptions}
\opt{A}{\bfigi{161.5}{TOAS01_p1_chem_s1_q04_A.png}}
\opt{B}{\bfigi{149.4}{TOAS01_p1_chem_s1_q04_B.png}}
\opt{C}{\bfigi{147.2}{TOAS01_p1_chem_s1_q04_C.png}}
\opt{D}{\bfigi{131.8}{TOAS01_p1_chem_s1_q04_D.png}}
\end{qoptions}
\end{bankq}
\begin{bankq}[id=TOAS01-P2-CHEM-S1-Q14, source={TOAS - Test 01 (30 Jan 2022)}, paper=2, qno=14, type=single, subtopic={Optical}, answer={A}, answer_src=key, difficulty=2, diff_parts={K1 S1 I1 C0 F0}, diff_note={concepts: enantiomers | steps: assign configuration of each pair; check mirror relationship | I: spot the pair that is superimposable}]
Which of the following pair of compounds doesn't represent enantiomeric pair :
\begin{qoptions}
\opt{A}{\bfigi{114.2}{TOAS01_p2_chem_s1_q14_A.png}}
\opt{B}{\bfigi{105.1}{TOAS01_p2_chem_s1_q14_B.png}}
\opt{C}{\bfigi{114.6}{TOAS01_p2_chem_s1_q14_C.png}}
\opt{D}{\bfigi{133.9}{TOAS01_p2_chem_s1_q14_D.png}}
\end{qoptions}
\end{bankq}
\begin{bankq}[id=TOAS01-P2-CHEM-S1-Q16, source={TOAS - Test 01 (30 Jan 2022)}, paper=2, qno=16, type=single, subtopic={Optical}, answer={D}, answer_src=key, difficulty=2, diff_parts={K1 S2 I0 C0 F0}, diff_note={concepts: symmetry elements | steps: CH2Cl2 has no inversion centre; second structure; cis-but-2-ene; last structure with both elements}]
Which of the following compound(s) contains plane of symmetry as well as centre of symmetry
\begin{qoptions}
\opt{A}{$\mathrm{CH_2Cl_2}$}
\opt{B}{\bfigi{58.0}{TOAS01_p2_chem_s1_q16_B.png}}
\opt{C}{Cis-but-2-ene}
\opt{D}{\bfigi{43.0}{TOAS01_p2_chem_s1_q16_2.png}}
\end{qoptions}
\end{bankq}
\begin{bankq}[id=TOAS01-P2-CHEM-S1-Q17, source={TOAS - Test 01 (30 Jan 2022)}, paper=2, qno=17, type=single, subtopic={Optical}, answer={B}, answer_src=key, difficulty=2, diff_parts={K1 S2 I0 C0 F0}, diff_note={concepts: meso compounds | steps: check each compound for symmetric stereocentres; check internal plane; find the one that cannot be meso}]
Which of the following compound(s) can not exist in meso form ?
\begin{qoptions}
\opt{A}{\bfigi{58.8}{TOAS01_p2_chem_s1_q17_A.png}}
\opt{B}{\bfigi{27.4}{TOAS01_p2_chem_s1_q17_B.png}}
\opt{C}{\bfigi{49.2}{TOAS01_p2_chem_s1_q17_C.png}}
\opt{D}{\bfigi{55.1}{TOAS01_p2_chem_s1_q17_D.png}}
\end{qoptions}
\end{bankq}
\begin{bankq}[id=TOAS01-P2-CHEM-S2-Q08, source={TOAS - Test 01 (30 Jan 2022)}, paper=2, qno=8, type=numerical, subtopic={Optical}, answer={2}, answer_src=solved, working={Chiral with $C_2$: the 2-ethyl-6-methyl biphenyl (atropisomer) and 1,3-dichloroallene; the others are achiral}, difficulty=4, diff_parts={K1 S2 I1 C0 F1}, diff_note={concepts: chirality and C2 axis | steps: identify chiral molecules; test for a C2 axis in each; count | I: chirality is compatible with C2 axis}]
Among the following, the number of molecules which are optically active \& possess $C_2$ axis of symmetry are :
\bfig{144.9}{TOAS01_p2_chem_s2_q08_1.png}
$\mathrm{CHCl_3}$ , 2,5-dimethylthiophene,
\bfig{62.0}{TOAS01_p2_chem_s2_q08_2.png}
\end{bankq}
\begin{bankq}[id=TOAS01-P2-CHEM-S2-Q09, source={TOAS - Test 01 (30 Jan 2022)}, paper=2, qno=9, type=numerical, subtopic={Geometrical}, answer={3}, answer_src=solved, working={Geometrical isomerism: but-2-ene, hexa-2,3,4-triene (odd cumulene), pent-2-ene; the spiro compound and the cyclobutylidene allenes show axial chirality, 2-ethylbut-1-ene none}, difficulty=2, diff_parts={K1 S1 I0 C0 F1}, diff_note={concepts: geometrical isomerism | steps: check each compound for a double bond or ring with two different substituents on each side}]
How many compound(s) can show geometrical isomerism ?
\bfig{150.1}{TOAS01_p2_chem_s2_q09_1.png}
\end{bankq}
\begin{bankq}[id=TOAS02-P1-CHEM-S1-Q03, source={TOAS - Test 02 (13 Feb 2022)}, paper=1, qno=3, type=single, subtopic={Optical}, answer={B}, answer_src=key, difficulty=1, diff_parts={K1 S1 I0 C0 F0}, diff_note={concepts: chirality and symmetry | steps: chirality of each structure; presence of a C2 axis}]
Which of the following compound is optically active with $\mathrm{C_2}$-axis of symmetry ?
\begin{qoptions}
\opt{A}{\bfigi{45.6}{TOAS02_p1_chem_s1_q03_A.png}}
\opt{B}{\bfigi{44.8}{TOAS02_p1_chem_s1_q03_B.png}}
\opt{C}{\bfigi{33.5}{TOAS02_p1_chem_s1_q03_C.png}}
\opt{D}{\bfigi{70.1}{TOAS02_p1_chem_s1_q03_D.png}}
\end{qoptions}
\end{bankq}
\begin{bankq}[id=TOAS02-P1-CHEM-S3-Q04, source={TOAS - Test 02 (13 Feb 2022)}, paper=1, qno=4, type=numerical, subtopic={Optical}, answer={5}, answer_src=solved, working={Chiral: (ii) 1,1,3-trimethylcyclopentane, (iii) amino alcohol, (v) 4-chlorocyclobutenone, (vi) spiro[3.3]heptane (axial), (vii) alkylidenecyclohexane (axial); (i) cumulene gives geometrical isomers, (iv) and (viii) achiral}, difficulty=4, diff_parts={K1 S2 I1 C0 F1}, diff_note={concepts: optical isomerism including axial chirality | steps: centres of chirality in each; axial chirality in allenes and spiro compounds; count | I: axial chirality}]
Number of molecules show optical isomerism :
\bfig{260.6}{TOAS02_p1_chem_s3_q04_1.png}
\end{bankq}
\begin{bankq}[id=TOAS03-P1-CHEM-S1-Q01, source={TOAS - Test 03 (27 Feb 2022)}, paper=1, qno=1, type=single, subtopic={Optical}, answer={C}, answer_src=key, difficulty=1, diff_parts={K1 S1 I0 C0 F0}, diff_note={concepts: resolution of optical isomers | steps: symmetry of options A and B; of C and D}]
Which of the following molecule is resolvable / optically active?
\begin{qoptions}
\opt{A}{\bfigi{41.3}{TOAS03_p1_chem_s1_q01_A.png}}
\opt{B}{\bfigi{81.1}{TOAS03_p1_chem_s1_q01_B.png}}
\opt{C}{\bfigi{53.2}{TOAS03_p1_chem_s1_q01_C.png}}
\opt{D}{\bfigi{123.0}{TOAS03_p1_chem_s1_q01_D.png}}
\end{qoptions}
\end{bankq}
\begin{bankq}[id=TOAS03-P1-CHEM-S3-Q01, source={TOAS - Test 03 (27 Feb 2022)}, paper=1, qno=1, type=numerical, subtopic={Optical}, answer={7}, answer_src=key, difficulty=3, diff_parts={K1 S2 I0 C0 F1}, diff_note={concepts: chiral centres in a polycyclic compound | steps: locate sp3 carbons; check four different groups for each; count 7}]
How many chiral center are present in this compound
\bfig{105.3}{TOAS03_p1_chem_s3_q01_1.png}
\end{bankq}
\begin{bankq}[id=TOAS03-P2-CHEM-S2-Q02, source={TOAS - Test 03 (27 Feb 2022)}, paper=2, qno=2, type=numerical, subtopic={Geometrical}, answer={6}, answer_src=solved, working={Stereo elements: S of sulfoxide (1), C=C (1), CHCl (1), ring C1 and C4 (2), C=N-OH (1): $2^6$ so $n=6$}, difficulty=5, diff_parts={K1 S3 I1 C0 F1}, diff_note={concepts: stereoisomers of a molecule with several elements | steps: stereocentres; C=C isomerism; chiral sulfur; oxime; ring centres | I: sulfoxide and oxime as stereo elements}]
\bfigi{198.5}{TOAS03_p2_chem_s2_q02_1.png}

If total number of possible stereoisomers for above molecule are $2^n$. The value of $n$ will be :
\end{bankq}
\begin{bankq}[id=TOAS05-P1-CHEM-S1-Q01, source={TOAS - Test 05 (03 Apr 2022)}, paper=1, qno=1, type=multi, subtopic={Geometrical}, answer={B,C,D}, answer_src=key, difficulty=3, diff_parts={K1 S2 I0 C0 F1}, diff_note={concepts: equilibrium constants of isomer pairs | steps: stability of each pair of isomers; direction of equilibrium; compare K}]
Which of the following is/are correctly match with their equilibrium constant ($K_{\mathrm{eq}}$)
\begin{qoptions}
\opt{A}{\bfigi{131.8}{TOAS05_p1_chem_s1_q01_A.png}}
\opt{B}{\bfigi{139.6}{TOAS05_p1_chem_s1_q01_B.png}}
\opt{C}{\bfigi{143.0}{TOAS05_p1_chem_s1_q01_C.png}}
\opt{D}{\bfigi{146.2}{TOAS05_p1_chem_s1_q01_D.png}}
\end{qoptions}
\end{bankq}
\begin{bankq}[id=TOAS05-P1-CHEM-S1-Q03, source={TOAS - Test 05 (03 Apr 2022)}, paper=1, qno=3, type=multi, subtopic={Optical}, answer={A,C}, answer_src=key, difficulty=2, diff_parts={K1 S1 I0 C0 F1}, diff_note={concepts: isomer relations | steps: identify A and B using the bromine water test; relation between them}]
\bfigi{58.0}{TOAS05_p1_chem_s1_q03_1.png} $A+B$\par
Identify relation between $A$ and $B$. If both $A$ and $B$ give Bromine water test
\begin{qoptions}
\opt{A}{Stereoisomers}
\opt{B}{Structural isomers}
\opt{C}{Enantiomers}
\opt{D}{Diastereomers}
\end{qoptions}
\end{bankq}
\begin{bankq}[id=TOAS05-P2-CHEM-S1-Q07, source={TOAS - Test 05 (03 Apr 2022)}, paper=2, qno=7, type=single, subtopic={Optical}, answer={B}, answer_src=key, difficulty=2, diff_parts={K1 S2 I0 C0 F0}, diff_note={concepts: stereochemical relation of products | steps: products B and D; chirality of each; relationship}]
$\mathrm{(1)\ CH_3OH}\xrightarrow{\mathrm{Na}}A$
\bfig{129.3}{TOAS05_p2_chem_s1_q06_2.png}
Product $B$ and $D$ are :
\begin{qoptions}
\opt{A}{Diastereomers}
\opt{B}{Enantiomers}
\opt{C}{Identical}
\opt{D}{Functional group isomers}
\end{qoptions}
\end{bankq}
\begin{bankq}[id=TOAS06-P1-CHEM-S2-Q08, source={TOAS - Test 06 (17 Apr 2022)}, paper=1, qno=8, type=multi, subtopic={Optical}, answer={A,C,D}, answer_src=key, difficulty=3, diff_parts={K1 S2 I0 C0 F1}, diff_note={concepts: symmetry and isomerism of a compound | steps: geometrical isomerism; centre of symmetry; plane of symmetry; optical isomerism}]
Correct statement(s) about this compound :
\bfig{69.2}{TOAS06_p1_chem_s2_q08_1.png}
\begin{qoptions}
\opt{A}{It shows geometrical isomerism}
\opt{B}{It has centre of symmetry}
\opt{C}{It has plane of symmetry}
\opt{D}{It shows optical isomerism}
\end{qoptions}
\end{bankq}
\begin{bankq}[id=TOAS16-P2-CHEM-S4-Q08, source={TOAS - Test 16 (04 Sep 2022)}, paper=2, qno=8, type=numerical, subtopic={Structural}, answer={3}, answer_src=key, difficulty=2, diff_parts={K1 S1 I0 C0 F1}, diff_note={concepts: isomers of C8H10 | steps: dimethylbenzenes give diacids; ethylbenzene does not}]
How many isomers ``X'' of C$_8$H$_{10}$ when reacts with hot alkaline KMnO$_4$ give only aromatic dicarboxylic acid?
\end{bankq}
\begin{bankq}[id=TOAS19-P1-CHEM-S1-Q05, source={TOAS - Test 19 (05 Feb 2023)}, paper=1, qno=5, type=multi, subtopic={Geometrical}, answer={B,C}, answer_src=key, difficulty=2, diff_parts={K1 S1 I0 C0 F1}, diff_note={concepts: stereochemistry of decalin | steps: identify the trans fusion from the drawings}, also_in={TNAS - Test 01 (05 Feb 2023) P2 Q3}]
Identify the structure of transdecaline
\begin{qoptions}
\opt{A}{\bfigi{89.7}{TOAS19_p1_chem_s1_q05_A.png}}
\opt{B}{\bfigi{112.3}{TOAS19_p1_chem_s1_q05_B.png}}
\opt{C}{\bfigi{79.4}{TOAS19_p1_chem_s1_q05_C.png}}
\opt{D}{\bfigi{71.3}{TOAS19_p1_chem_s1_q05_D.png}}
\end{qoptions}
\end{bankq}
\begin{bankq}[id=TOAS19-P1-CHEM-S1-Q06, source={TOAS - Test 19 (05 Feb 2023)}, paper=1, qno=6, type=multi, subtopic={Optical}, answer={B,C}, answer_src=key, difficulty=3, diff_parts={K1 S2 I0 C0 F1}, diff_note={concepts: identity of stereoisomeric drawings | steps: compare configurations; symmetry; identical pairs}, also_in={TNAS - Test 01 (05 Feb 2023) P2 Q4}]
Identify the structure of identical compounds
\begin{qoptions}
\opt{A}{\bfigi{65.6}{TOAS19_p1_chem_s1_q06_A.png}}
\opt{B}{\bfigi{70.3}{TOAS19_p1_chem_s1_q06_B.png}}
\opt{C}{\bfigi{70.8}{TOAS19_p1_chem_s1_q06_C.png}}
\opt{D}{\bfigi{70.3}{TOAS19_p1_chem_s1_q06_D.png}}
\end{qoptions}
\end{bankq}
\begin{bankq}[id=TOAS19-P1-CHEM-S1-Q07, source={TOAS - Test 19 (05 Feb 2023)}, paper=1, qno=7, type=multi, subtopic={Geometrical}, answer={A,C}, answer_src=key, difficulty=2, diff_parts={K1 S1 I0 C0 F1}, diff_note={concepts: geometrical isomerism | steps: identify cis and trans arrangements; compare each pair}, also_in={TNAS - Test 01 (05 Feb 2023) P2 Q5}]
Identify the geometrical isomers
\begin{qoptions}
\opt{A}{\bfigi{95.1}{TOAS19_p1_chem_s1_q07_A.png}}
\opt{B}{\bfigi{87.4}{TOAS19_p1_chem_s1_q07_B.png}}
\opt{C}{\bfigi{89.7}{TOAS19_p1_chem_s1_q07_C.png}}
\opt{D}{\bfigi{92.5}{TOAS19_p1_chem_s1_q07_D.png}}
\end{qoptions}
\end{bankq}

\section*{TOPS papers}
\begin{bankpassage}[id=TOPS04-P1-CHEM-S3-G3-4, source={TOPS - Test 04 (18 Dec 2022)}, questions={3-4}]
Compounds having same molecular formula same structure but different 3D arrangement of atoms or groups are called stereoisomer of each other and phenomenon is called stereo isomerism.
\end{bankpassage}
\begin{bankq}[id=TOPS04-P1-CHEM-S3-Q03, source={TOPS - Test 04 (18 Dec 2022)}, paper=1, qno=3, type=numerical, subtopic={Geometrical}, answer={5.00}, answer_src=key, difficulty=5, diff_parts={K1 S2 I1 C1 F1}, diff_note={concepts: geometrical isomerism | steps: identify double bonds and rings; end groups; count}, group=TOPS04-P1-CHEM-S3-G3-4]
Possible number of G.I. of compound?
\bfig{68.6}{TOPS04_p1_chem_s3_q03_1.png}
\end{bankq}
\begin{bankq}[id=TOPS04-P1-CHEM-S3-Q04, source={TOPS - Test 04 (18 Dec 2022)}, paper=1, qno=4, type=numerical, subtopic={Geometrical}, answer={0.00}, answer_src=key, difficulty=5, diff_parts={K1 S2 I1 C1 F1}, diff_note={concepts: chirality | steps: identify each carbon; four different groups; count}, group=TOPS04-P1-CHEM-S3-G3-4]
How many chiral carbon are present in
\bfig{46.0}{TOPS04_p1_chem_s3_q04_1.png}
\end{bankq}
\begin{bankq}[id=TOPS04-P2-CHEM-S1-Q13, source={TOPS - Test 04 (18 Dec 2022)}, paper=2, qno=13, type=single, subtopic={Optical}, answer={C}, answer_src=key, difficulty=3, diff_parts={K1 S1 I1 C1 F0}, diff_note={concepts: stereochemical relation | steps: configuration of each centre; compare}]
(A) \bfig{84.5}{TOPS04_p2_chem_s1_q13_1.png}
(B) \bfig{58.3}{TOPS04_p2_chem_s1_q13_B.png}
Find out the correct relation in the above pair of compound?
\begin{qoptions}
\opt{A}{Diastereomers}
\opt{B}{Enantiomers}
\opt{C}{Position isomers}
\opt{D}{Conformation}
\end{qoptions}
\end{bankq}
\begin{bankq}[id=TOPS04-P2-CHEM-S2-Q06, source={TOPS - Test 04 (18 Dec 2022)}, paper=2, qno=6, type=numerical, subtopic={Optical}, answer={4}, answer_src=key, difficulty=5, diff_parts={K1 S2 I1 C1 F1}, diff_note={concepts: chiral centres | steps: locate the carbons; four different groups; count}]
How many chiral centres are present in the following compound ?
\bfig{87.2}{TOPS04_p2_chem_s2_q06_1.png}
\end{bankq}
\begin{bankq}[id=TOPS04-P2-CHEM-S2-Q07, source={TOPS - Test 04 (18 Dec 2022)}, paper=2, qno=7, type=numerical, subtopic={Geometrical}, answer={4}, answer_src=solved, working={Tetrachlorocyclobutane with Cl on each carbon: up/down patterns UUUU, UUUD, UUDD and UDUD are distinct under the ring symmetry, so 4}, difficulty=5, diff_parts={K1 S2 I1 C1 F1}, diff_note={concepts: geometrical isomerism | steps: ring substituents; arrangement patterns; count}]
How many geometrical isomerse are possible in the given compound?
\bfig{44.8}{TOPS04_p2_chem_s2_q07_1.png}
\end{bankq}
\begin{bankq}[id=TOPS05-P1-CHEM-S1-Q03, source={TOPS - Test 05 (14 Jan 2023)}, paper=1, qno=3, type=multi, subtopic={Optical}, answer={A,B,C,D}, answer_src=key, difficulty=5, diff_parts={K1 S2 I1 C1 F1}, diff_note={concepts: compounds with a C3 axis | steps: determine the symmetry elements of each structure; C3 axis requires three-fold rotational symmetry; check four}]
Identify compounds having $C_3$ axis of symmetry?
\begin{qoptions}
\opt{A}{\bfigi{193.2}{TOPS05_p1_chem_s1_q03_A.png}}
\opt{B}{\bfigi{105.8}{TOPS05_p1_chem_s1_q03_B.png}\par\bfigi{95.8}{TOPS05_p1_chem_s1_q03_3.png}}
\opt{C}{\bfigi{172.9}{TOPS05_p1_chem_s1_q03_C.png}}
\opt{D}{\bfigi{154.3}{TOPS05_p1_chem_s1_q03_D.png}}
\end{qoptions}
\end{bankq}
\begin{bankq}[id=TOPS05-P1-CHEM-S1-Q04, source={TOPS - Test 05 (14 Jan 2023)}, paper=1, qno=4, type=multi, subtopic={Optical}, answer={A,B,C}, answer_src=key, difficulty=4, diff_parts={K1 S2 I1 C0 F1}, diff_note={concepts: species having C3 axis | steps: benzene-like ring rotation; inorganic benzene has C3; PCl5 has C3 axis; water has only C2}]
Identify compounds having $C_3$ axis of symmetry?
\begin{qoptions}
\opt{A}{\bfigi{44.1}{TOPS05_p1_chem_s1_q04_A.png}}
\opt{B}{Inorganic benzene}
\opt{C}{$\mathrm{PCl_5}$}
\opt{D}{Water}
\end{qoptions}
\end{bankq}
\begin{bankq}[id=TOPS05-P1-CHEM-S2-Q09, source={TOPS - Test 05 (14 Jan 2023)}, paper=1, qno=9, type=single, subtopic={Optical}, answer={A}, answer_src=key, difficulty=4, diff_parts={K1 S2 I1 C1 F0}, diff_note={concepts: optical activity of drawn compounds | steps: check chirality or meso nature for each structure; plane of symmetry; match}]
{\footnotesize\setlength{\tabcolsep}{2pt}
\begin{tabular}{|l|p{3.1cm}|l|p{3.0cm}|}
\hline
 & List-I & & List-II\\
\hline
(P) & \bfigi{50.4}{TOPS05_p1_chem_s2_q09_1.png} & (1) & Optically inactive compound\\
(Q) & \bfigi{105.6}{TOPS05_p1_chem_s2_q09_2.png} & (2) & Optically active compound\\
(R) & \bfigi{50.4}{TOPS05_p1_chem_s2_q09_3.png} & (3) & Dissymmetric compound\\
(S) & \bfigi{67.1}{TOPS05_p1_chem_s2_q09_4.png} & (4) & Axis of symmetry present\\
 & & (5) & Asymmetric compound\\
\hline
\end{tabular}}
\begin{qoptions}
\opt{A}{P$\to$1;Q$\to$1,4;R$\to$2,5;S$\to$2,5}
\opt{B}{P$\to$2;Q$\to$3,4;R$\to$1,5;S$\to$2,5}
\opt{C}{P$\to$1,2;Q$\to$3,4;R$\to$4,5;S$\to$4}
\opt{D}{P$\to$3;Q$\to$2;R$\to$1;S$\to$1,2}
\end{qoptions}
\end{bankq}
\begin{bankq}[id=TOPS05-P1-CHEM-S3-Q03, source={TOPS - Test 05 (14 Jan 2023)}, paper=1, qno=3, type=numerical, subtopic={Geometrical}, answer={8.00}, answer_src=key, difficulty=5, diff_parts={K1 S2 I1 C1 F1}, diff_note={concepts: total number of stereoisomers of a drawn compound | steps: identify stereocentres and double bonds; check for meso forms and symmetry; count}]
Identify total number of stereo isomers for following structure
\bfig{76.8}{TOPS05_p1_chem_s3_q03_1.png}
\end{bankq}
\begin{bankq}[id=TOPS05-P1-CHEM-S3-Q05, source={TOPS - Test 05 (14 Jan 2023)}, paper=1, qno=5, type=numerical, subtopic={Optical}, answer={4.00}, answer_src=key, difficulty=5, diff_parts={K1 S2 I1 C1 F1}, diff_note={concepts: total stereoisomers of a compound | steps: count chiral centres and geometric isomerism; account for symmetry; total}]
\bfig{108.2}{TOPS05_p1_chem_s3_q05_1.png}
Total number of stereoisomers in the above compound is?
\end{bankq}
\begin{bankq}[id=TOPS05-P2-CHEM-S2-Q08, source={TOPS - Test 05 (14 Jan 2023)}, paper=2, qno=8, type=multi, subtopic={Optical}, answer={C,D}, answer_src=key, difficulty=4, diff_parts={K1 S2 I1 C0 F1}, diff_note={concepts: enantiomeric pairs among structures | steps: assign configurations at each stereocentre; mirror image relationship; check four}]
Among the following pair(s) of compounds identify enantiomeric pairs
\begin{qoptions}
\opt{A}{\bfigi{108.0}{TOPS05_p2_chem_s2_q08_A.png}}
\opt{B}{\bfigi{120.5}{TOPS05_p2_chem_s2_q08_B.png}}
\opt{C}{\bfigi{107.8}{TOPS05_p2_chem_s2_q08_C.png}}
\opt{D}{\bfigi{55.4}{TOPS05_p2_chem_s2_q08_D.png}\ \bfigi{42.3}{TOPS05_p2_chem_s2_q08_5.png}}
\end{qoptions}
\end{bankq}
\begin{bankq}[id=TOPS05-P2-CHEM-S2-Q09, source={TOPS - Test 05 (14 Jan 2023)}, paper=2, qno=9, type=multi, subtopic={Optical}, answer={B}, answer_src=key, difficulty=5, diff_parts={K1 S2 I1 C1 F1}, diff_note={concepts: specific rotation from observed rotation | steps: concentration g/mL from 24 g in 240 mL; path length in dm; specific rotation = alpha/(c l); evaluate}]
24 gm of optically pure tartaric acid is dissolved in water to make 240 mL solution. It is kept in 20 cm polarimeter tube and plane polarised light is passed through it to produce rotation of $-2.4^\circ$. Calculate specific rotation for above compound?
\begin{qoptions}
\opt{A}{$12^\circ$}
\opt{B}{$-12^\circ$}
\opt{C}{$+6^\circ$}
\opt{D}{$-6^\circ$}
\end{qoptions}
\end{bankq}
\begin{bankq}[id=TOPS05-P2-CHEM-S3-Q05, source={TOPS - Test 05 (14 Jan 2023)}, paper=2, qno=5, type=numerical, subtopic={Optical}, answer={3}, answer_src=key, difficulty=4, diff_parts={K1 S1 I1 C1 F1}, diff_note={concepts: chiral centres in a drawn compound | steps: identify sp3 carbons with four different groups; count}]
Identify total number of chiral centers in following compound?
\bfig{130.8}{TOPS05_p2_chem_s3_q05_1.png}
\end{bankq}

\section*{TEAS papers}
\begin{bankq}[id=TEAS01-P1-CHEM-S1-Q04, source={TEAS - Test 01 (16 Apr 2023)}, paper=1, qno=4, type=multi, subtopic={Optical}, answer={A,B,D}, answer_src=key, difficulty=6, diff_parts={K1 S2 I2 C1 F1}, diff_note={concepts: stereochemistry of reactions | steps: stereocentres formed; relation of the products; check each}]
Which of the following reaction result the formation of a pair of diastereomers?
\begin{qoptions}
\opt{A}{\bfigi{92.5}{TEAS01_p1_chem_s1_q04_A.png}}
\opt{B}{\bfigi{78.2}{TEAS01_p1_chem_s1_q04_B.png}}
\opt{C}{\bfigi{90.9}{TEAS01_p1_chem_s1_q04_C.png}}
\opt{D}{\bfigi{34.2}{TEAS01_p1_chem_s1_q04_D.png}}
\end{qoptions}
\end{bankq}

\section*{TMAS papers}
\begin{bankq}[id=TMAS01-P1-CHEM-S1-Q15, source={TMAS - Test 01 (09 Jul 2023)}, paper=1, qno=15, type=single, subtopic={Optical}, answer={B}, answer_src=key, difficulty=1, diff_parts={K1 S1 I0 C0 F0}, diff_note={concepts: stereoisomerism | steps: configuration at each centre; relationship of the products}]
\par
\bfig{165.1}{TMAS01_p1_chem_s1_q15_1.png}
Identify relationship between produts ?
\begin{qoptions}
\opt{A}{Diastereo Isomers}
\opt{B}{Enantiomers}
\opt{C}{Functional Isomers}
\opt{D}{Geometrical Isomers}
\end{qoptions}
\end{bankq}
\begin{bankq}[id=TMAS02-P2-CHEM-S1-Q05, source={TMAS - Test 02 (06 Aug 2023)}, paper=2, qno=5, type=multi, subtopic={Optical}, answer={A,D}, answer_src=solved, working={Syn dihydroxylation of trans-2-butene gives a racemic pair; epoxidation then anti opening gives the meso diol}, difficulty=6, diff_parts={K1 S2 I2 C1 F1}, diff_note={concepts: stereochemistry of addition | steps: syn and anti addition; products; meso or racemic}]
\bfig{172.3}{TMAS02_p2_chem_s1_q05_1.png}
Select correct statement(s) for above reaction sequence :
\begin{qoptions}
\opt{A}{$A$ is pair of enantiomer and $A$ and $B$ are distereomers}
\opt{B}{$B$ is pair of enantiomer and $A$ and $B$ are distereomers}
\opt{C}{$A$ represents meso compound}
\opt{D}{$B$ represents meso compound}
\end{qoptions}
\end{bankq}

\section*{TNAS papers}
\begin{bankq}[id=TNAS01-P1-CHEM-S1-Q19, source={TNAS - Test 01 (05 Feb 2023)}, paper=1, qno=19, type=single, subtopic={Optical}, answer={A}, answer_src=key, difficulty=1, diff_parts={K1 S1 I0 C0 F0}, diff_note={concepts: stereoisomers | steps: assign configurations at each centre; check which pair differs at only some centres}]
Identify Diastereoisomeric pair of compounds?
\begin{qoptions}
\opt{A}{\bfigi{125.8}{TNAS01_p1_chem_s1_q19_A.png}}
\opt{B}{\bfigi{135.2}{TNAS01_p1_chem_s1_q19_B.png}}
\opt{C}{\bfigi{128.3}{TNAS01_p1_chem_s1_q19_C.png}}
\opt{D}{\bfigi{133.7}{TNAS01_p1_chem_s1_q19_D.png}}
\end{qoptions}
\end{bankq}
\begin{bankq}[id=TNAS01-P1-CHEM-S2-Q08, source={TNAS - Test 01 (05 Feb 2023)}, paper=1, qno=8, type=numerical, subtopic={Optical}, answer={3}, answer_src=solved, working={Chiral: sec-butanol, the allene with different ends, and trans-1,3-dichlorocyclopentane; the norbornene, spiro, cis-cyclopropane and cis-dioxetane compounds are achiral. Count 3}, difficulty=2, diff_parts={K1 S1 I0 C0 F1}, diff_note={concepts: optical activity | steps: identify stereocentres and symmetry elements in each compound; count chiral}]
Identify total number of compounds which are optically active?
\bfig{128.6}{TNAS01_p1_chem_s2_q08_1.png}
\bfig{140.4}{TNAS01_p1_chem_s2_q08_2.png}
\bfig{131.5}{TNAS01_p1_chem_s2_q08_3.png}
\bfig{74.5}{TNAS01_p1_chem_s2_q08_4.png}
\end{bankq}
\begin{bankq}[id=TNAS01-P1-CHEM-S2-Q10, source={TNAS - Test 01 (05 Feb 2023)}, paper=1, qno=10, type=numerical, subtopic={Geometrical}, answer={7}, answer_src=solved, working={Four stereocentres CH(OH)CH3 on a tetrasubstituted alkene; counting R/S labelings of the four in-plane positions modulo the proper symmetries gives $(16+4+4+4)/4=7$}, difficulty=5, diff_parts={K1 S2 I1 C1 F1}, diff_note={concepts: stereoisomer count | steps: count stereocentres and double bonds; symmetry reduces; total stereoisomers | I: symmetry halves some meso combinations}]
Total number of stereo isomers for following structure?
\bfig{50.2}{TNAS01_p1_chem_s2_q10_1.png}
\end{bankq}
\begin{bankq}[id=TNAS03-P1-CHEM-S1-Q19, source={TNAS - Test 03 (28 May 2023)}, paper=1, qno=19, type=single, subtopic={Optical}, answer={C}, answer_src=key, difficulty=3, diff_parts={K1 S1 I1 C1 F0}, diff_note={concepts: enantiomer in Fischer projection | steps: configuration; mirror image}]
Which of the following Fischer projection is the enantiomer of the following molecule?
\bfig{87.4}{TNAS03_p1_chem_s1_q19_1.png}
\begin{qoptions}
\opt{A}{\bfigi{61.5}{TNAS03_p1_chem_s1_q19_A.png}}
\opt{B}{\bfigi{58.3}{TNAS03_p1_chem_s1_q19_B.png}}
\opt{C}{\bfigi{61.5}{TNAS03_p1_chem_s1_q19_C.png}}
\opt{D}{\bfigi{58.3}{TNAS03_p1_chem_s1_q19_D.png}}
\end{qoptions}
\end{bankq}
\begin{bankq}[id=TNAS03-P1-CHEM-S2-Q09, source={TNAS - Test 03 (28 May 2023)}, paper=1, qno=9, type=numerical, subtopic={Optical}, answer={1}, answer_src=solved, working={3,4-Dibromo ring with a symmetric allene end gives one meso (cis) form; the trans form is chiral}, difficulty=6, diff_parts={K1 S2 I2 C1 F1}, diff_note={concepts: meso isomers | steps: stereocentres and axes; symmetry; count}]
Number of meso isomer of the given compound is :
\bfig{50.4}{TNAS03_p1_chem_s2_q09_1.png}
\end{bankq}
\begin{bankq}[id=TNAS03-P2-CHEM-S1-Q04, source={TNAS - Test 03 (28 May 2023)}, paper=2, qno=4, type=multi, subtopic={Optical}, answer={A,B,C,D}, answer_src=key, difficulty=6, diff_parts={K1 S2 I2 C1 F1}, diff_note={concepts: stereochemistry of a sequence | steps: products; configuration; planarity; check}]
Consider following reactions
\bfig{138.5}{TNAS03_p2_chem_s1_q04_1.png}
\bfig{127.3}{TNAS03_p2_chem_s1_q04_2.png}
Choose incorrect options :
\begin{qoptions}
\opt{A}{Total number of maximum atom in one plane is 8 in product B}
\opt{B}{D and F are identical}
\opt{C}{Q to E is formed by $\mathrm{S_N2}$ mechanism}
\opt{D}{Configuration is oppossite in D and F but same in D and B at chiral centre}
\end{qoptions}
\end{bankq}
\begin{bankq}[id=TNAS04-P2-CHEM-S3-Q06, source={TNAS - Test 04 (18 Jun 2023)}, paper=2, qno=6, type=numerical, subtopic={Optical}, answer={4.00}, answer_src=key, difficulty=5, diff_parts={K1 S2 I1 C1 F1}, diff_note={concepts: stereoisomers of an unsaturated alcohol | steps: chiral centre; double bond geometry; combine; count}]
$\mathrm{Me-\overset{OH}{\overset{|}{CH}}-CH{=}CH-Me}$\par
Total number of stereoisomers?
\end{bankq}
\begin{bankq}[id=TNAS05-P1-CHEM-S1-Q16, source={TNAS - Test 05 (09 Jul 2023)}, paper=1, qno=16, type=single, subtopic={Optical}, answer={D}, answer_src=key, difficulty=3, diff_parts={K1 S2 I1 C0 F0}, diff_note={concepts: retention of optical activity | steps: check mechanism of each reaction; SN1 gives racemic; the one that keeps the stereocentre | I: stereocentre not involved in the reaction}]
In which of the following reaction the product will also be optically active?
\begin{qoptions}
\opt{A}{\bfigi{133.0}{TNAS05_p1_chem_s1_q16_A.png}}
\opt{B}{\bfigi{77.9}{TNAS05_p1_chem_s1_q16_B.png}}
\opt{C}{\bfigi{114.1}{TNAS05_p1_chem_s1_q16_C.png}}
\opt{D}{\bfigi{89.4}{TNAS05_p1_chem_s1_q16_D.png}}
\end{qoptions}
\end{bankq}
\begin{bankq}[id=TNAS05-P1-CHEM-S1-Q18, source={TNAS - Test 05 (09 Jul 2023)}, paper=1, qno=18, type=single, subtopic={Optical}, answer={C}, answer_src=key, difficulty=3, diff_parts={K1 S2 I1 C0 F0}, diff_note={concepts: reactivity of an alkene and chirality | steps: optical activity of P; HBr addition product; bromine water test; fourth statement}]
\bfig{95.0}{TNAS05_p1_chem_s1_q18_1.png}
Choose correct option?
\begin{qoptions}
\opt{A}{P is optically active.}
\opt{B}{Reaction of P with HBr give optically active product.}
\opt{C}{P can give bromine water test.}
\opt{D}{Reaction of P with $\mathrm{HBr/R_2O_2}$ give one chiral centre containing product.}
\end{qoptions}
\end{bankq}

\section*{Star Review (Other) papers}
\begin{bankq}[id=STAR03-P1-CHEM-S2-Q07, source={Star Review (Other) - Test 03 (04 Dec 2023)}, paper=1, qno=7, type=multi, subtopic={Optical}, answer={C,D}, answer_src=key, difficulty=6, diff_parts={K1 S2 I2 C1 F1}, diff_note={concepts: specific rotation of products of nitration | steps: identify the structures from the data; products from each; optical rotation; match}]
Given :\par
\bfig{165.3}{STAR03_p1_chem_s2_q07_1.png}
The compound(s), which on reaction with $\mathrm{HNO_3}$ will give the product having degree of rotation $[\alpha]_D=-52.7^\circ$ is (are)
\begin{qoptions}
\opt{A}{\bfigi{49.9}{STAR03_p1_chem_s2_q07_A.png}}
\opt{B}{\bfigi{51.5}{STAR03_p1_chem_s2_q07_B.png}}
\opt{C}{\bfigi{51.6}{STAR03_p1_chem_s2_q07_C.png}}
\opt{D}{\bfigi{55.7}{STAR03_p1_chem_s2_q07_D.png}}
\end{qoptions}
\end{bankq}

\section*{Internal Test papers}
\begin{bankq}[id=INT02-P1-CHEM-S2-Q07, source={Internal Test - Test 02 (07 Jul 2021)}, paper=1, qno=7, type=multi, subtopic={Structural}, answer={A,C,D}, answer_src=key, difficulty=4, diff_parts={K1 S2 I1 C0 F1}, diff_note={concepts: keto-enol tautomerism | steps: check alpha hydrogens in each compound; the enol content; those that show tautomerism}]
Tautomerism is exhibited by :
\begin{qoptions}
\opt{A}{\bfigi{106.5}{INT02_p1_chem_s2_q07_A.png}}
\opt{B}{\bfigi{76.8}{INT02_p1_chem_s2_q07_B.png}}
\opt{C}{\bfigi{57.6}{INT02_p1_chem_s2_q07_C.png}}
\opt{D}{\bfigi{57.6}{INT02_p1_chem_s2_q07_D.png}}
\end{qoptions}
\end{bankq}
\begin{bankq}[id=INT02-P1-CHEM-S3-Q14, source={Internal Test - Test 02 (07 Jul 2021)}, paper=1, qno=14, type=numerical, subtopic={Optical}, answer={1}, answer_src=key, difficulty=2, diff_parts={K1 S1 I0 C0 F1}, diff_note={concepts: symmetry elements | steps: locate any plane of symmetry in each molecule; count those without}]
How many of the following do not have plane of symmetry ?\par
\bfig{214.9}{INT02_p1_chem_s3_q14_1.png}
\bfig{97.7}{INT02_p1_chem_s3_q14_2.png}
\bfig{58.5}{INT02_p1_chem_s3_q14_3.png}
\end{bankq}
\end{document}
